% Options for packages loaded elsewhere
\PassOptionsToPackage{unicode}{hyperref}
\PassOptionsToPackage{hyphens}{url}
\PassOptionsToPackage{dvipsnames,svgnames,x11names}{xcolor}
\documentclass[
  11pt,
  a4paper,
  11pt,
  oneside]{report}
\usepackage{xcolor}
\usepackage[margin=1in,top=1.25in,bottom=1.5in]{geometry}
\usepackage{amsmath,amssymb}
\setcounter{secnumdepth}{-\maxdimen} % remove section numbering
\usepackage{iftex}
\ifPDFTeX
  \usepackage[T1]{fontenc}
  \usepackage[utf8]{inputenc}
  \usepackage{textcomp} % provide euro and other symbols
\else % if luatex or xetex
  \usepackage{unicode-math} % this also loads fontspec
  \defaultfontfeatures{Scale=MatchLowercase}
  \defaultfontfeatures[\rmfamily]{Ligatures=TeX,Scale=1}
\fi
\usepackage{lmodern}
\ifPDFTeX\else
  % xetex/luatex font selection
  \setmainfont[Ligatures=TeX]{Inter}
  \setmonofont[Scale=0.91]{JetBrains Mono}
\fi
% Use upquote if available, for straight quotes in verbatim environments
\IfFileExists{upquote.sty}{\usepackage{upquote}}{}
\IfFileExists{microtype.sty}{% use microtype if available
  \usepackage[]{microtype}
  \UseMicrotypeSet[protrusion]{basicmath} % disable protrusion for tt fonts
}{}
\usepackage{setspace}
\makeatletter
\@ifundefined{KOMAClassName}{% if non-KOMA class
  \IfFileExists{parskip.sty}{%
    \usepackage{parskip}
  }{% else
    \setlength{\parindent}{0pt}
    \setlength{\parskip}{6pt plus 2pt minus 1pt}}
}{% if KOMA class
  \KOMAoptions{parskip=half}}
\makeatother
\usepackage{longtable,booktabs,array}
\usepackage{calc} % for calculating minipage widths
% Correct order of tables after \paragraph or \subparagraph
\usepackage{etoolbox}
\makeatletter
\patchcmd\longtable{\par}{\if@noskipsec\mbox{}\fi\par}{}{}
\makeatother
% Allow footnotes in longtable head/foot
\IfFileExists{footnotehyper.sty}{\usepackage{footnotehyper}}{\usepackage{footnote}}
\makesavenoteenv{longtable}
\usepackage{graphicx}
\makeatletter
\newsavebox\pandoc@box
\newcommand*\pandocbounded[1]{% scales image to fit in text height/width
  \sbox\pandoc@box{#1}%
  \Gscale@div\@tempa{\textheight}{\dimexpr\ht\pandoc@box+\dp\pandoc@box\relax}%
  \Gscale@div\@tempb{\linewidth}{\wd\pandoc@box}%
  \ifdim\@tempb\p@<\@tempa\p@\let\@tempa\@tempb\fi% select the smaller of both
  \ifdim\@tempa\p@<\p@\scalebox{\@tempa}{\usebox\pandoc@box}%
  \else\usebox{\pandoc@box}%
  \fi%
}
% Set default figure placement to htbp
\def\fps@figure{htbp}
\makeatother
\setlength{\emergencystretch}{3em} % prevent overfull lines
\providecommand{\tightlist}{%
  \setlength{\itemsep}{0pt}\setlength{\parskip}{0pt}}
\usepackage{fancyhdr}
\usepackage{vcell}

% Custom colors
\usepackage[table]{xcolor}
\definecolor{primarycolor}{RGB}{13, 5, 35} % #0d0523
\definecolor{accentcolor}{RGB}{43, 55, 80} % #2b3750
\definecolor{sectioncolor}{RGB}{48, 26, 76} % #301a4c
\definecolor{emphasiscolor}{RGB}{105, 84, 150} % #695496
\definecolor{linkcolor}{RGB}{0, 128, 191} % #0080bf 
\definecolor{snippets}{RGB}{239, 242, 248} % #eff2f8



\renewcommand{\arraystretch}{1.3}

% Graphics for logo and background images
\usepackage{graphicx}
\usepackage{eso-pic}

% Prevent floats from overlapping headers/footers
\usepackage{float}

% Package for ensuring minimum space before blocks
\usepackage{needspace}


% Footer container to enable colored background
% Add shaded background to footer on every page
\AddToShipoutPictureBG{%
    \AtPageLowerLeft{%
        \color{primarycolor}%
        \rule{\paperwidth}{2.2cm}%  <--- Adjusted to match footer content height
    }%
}





% Prevent widows and orphans (single lines at top/bottom of pages)
\widowpenalty=10000
\clubpenalty=10000
\raggedbottom

% Prevent page breaks between paragraphs and following content
\predisplaypenalty=10000

% Increase penalty for content at bottom of page to prevent footer overlap
\makeatletter
\setlength{\@fpsep}{12pt plus 1fil}
\setlength{\@fpbot}{0pt plus 1fil}
\makeatother

% Control float placement to prevent overlapping with headers/footers
\setlength{\textfloatsep}{10pt}     % Minimum space between float and text
\setlength{\floatsep}{10pt}         % Space between consecutive floats
\renewcommand{\topfraction}{0.85}   % Max fraction of page for floats at top
\renewcommand{\bottomfraction}{0.7} % Max fraction of page for floats at bottom
\renewcommand{\textfraction}{0.15}  % Min fraction of page for text
\renewcommand{\floatpagefraction}{0.8} % Min fraction of float page that must be floats

% Reduce excessive spacing around figures (specific to figures only via AtBeginEnvironment below)
\setlength{\intextsep}{6pt}        % Space above and below in-text floats (figures)
\setlength{\abovecaptionskip}{4pt} % Space above caption
\setlength{\belowcaptionskip}{2pt} % Space below caption

% Apply tight spacing specifically to figure environments only (not tables)
\AtBeginEnvironment{figure}{\vspace{0pt}}
\AtEndEnvironment{figure}{\vspace{-4pt}}

% Title page - appears BEFORE TOC
\AtBeginDocument{%
\begin{titlepage}
\thispagestyle{empty}



% Background image - full page
\AddToShipoutPictureBG*{%
  
\includegraphics[width=\paperwidth,height=\paperheight]{images/brand-background.png}%
}

% Logo top right
\begin{flushright}
\vspace*{.5cm}

\includegraphics[width=0.1\textwidth]{images/logo.png}
\end{flushright}

\vspace*{2cm}

% Title centered (using Inter font)
\centering
{\Huge\bfseries\color{white} Visual Layer for Camtek}

\vspace{.8cm}

{\LARGE\bfseries\color{white} User Guide and Documentation}

\vspace{15cm}

% Auto-updating date (Month YYYY format) + Confidential Proprietary
{\normalsize\color{white} Confidential Proprietary | \ifcase\month\or January\or February\or March\or April\or May\or June\or July\or August\or September\or October\or November\or December\fi\space\number\year}

\vfill

\end{titlepage}

\newpage
\pagenumbering{arabic}
\setcounter{page}{1}
}

% Section counter for major sections
\newcounter{sectioncounter}

% Define custom section divider style with automatic numbering
\newcommand{\sectiondivider}[1]{%
  \stepcounter{sectioncounter}%
  \newpage
  \thispagestyle{empty}
  % Set page background
  \AddToShipoutPictureBG*{%
    \color[RGB]{242,244,248}\rule{\paperwidth}{\paperheight}%
  }
  \vspace*{\fill}
  \begin{center}
 \uppercase{\bfseries}{\Large\color{sectioncolor}SECTION \Roman{sectioncounter}}\\[2cm]
  {\Huge\color{sectioncolor}#1}
  \end{center}
  \vspace*{\fill}
  \newpage
  \addcontentsline{toc}{sectiondivider}{Section \Roman{sectioncounter}: #1}%
}

% Generic divider for Appendices (Same style as Section Divider)
\newcommand{\genericdivider}[1]{%
    \newpage
    \thispagestyle{empty}
    \AddToShipoutPictureBG*{%
        \color[RGB]{242,244,248}\rule{\paperwidth}{\paperheight}%
    }
    \vspace*{\fill}
    \begin{center}
        % No "SECTION X" counter here, just the title
        {\Huge\color{sectioncolor}#1}
    \end{center}
    \vspace*{\fill}
    \newpage
    \addcontentsline{toc}{genericdivider}{#1}%
}

% Appendix divider (styled like section dividers but unnumbered)
\newcommand{\appendixdivider}[1]{%
    \newpage
    \thispagestyle{empty}
    \AddToShipoutPictureBG*{%
        \color[RGB]{242,244,248}\rule{\paperwidth}{\paperheight}%
    }
    \vspace*{\fill}
    \begin{center}
        {\Huge\color{sectioncolor}#1}
    \end{center}
    \vspace*{\fill}
    \newpage
    \addcontentsline{toc}{appendixdivider}{#1}%
}

% Appendix counter and formatting (DEPRECATED - use \appendix instead)
\newcounter{appendixcounter}
\newcommand{\appendixheading}[1]{%
  \newpage
  \stepcounter{appendixcounter}%
  \section*{Appendix \Roman{appendixcounter}: #1}%
  \addcontentsline{toc}{section}{References \Roman{appendixcounter}: #1}%
}

% Configure appendix numbering to use Arabic numerals (1, 2, 3...)
% This overrides the default letter-based numbering (A, B, C...)
% For 'report' document class which uses chapters
\usepackage{titlesec}


\usepackage{tocloft}

\makeatletter
\addtocounter{tocdepth}{-2} 
% --- Chapter dots + spacing ---
\@ifundefined{cftchapdotsep}{}{\renewcommand{\cftchapdotsep}{\cftdotsep}}
\@ifundefined{cftchapleader}{}{\renewcommand{\cftchapleader}{\cftdotfill{\cftchapdotsep}}}
\@ifundefined{cftafterchapskip}{}{\setlength{\cftafterchapskip}{6pt}}
\@ifundefined{cftbeforechapskip}{}{\setlength{\cftbeforechapskip}{8pt}}

% Ensure chapter TOC entries always have dotted leaders and spacing
\renewcommand{\l@chapter}[2]{%
  \@dottedtocline{0}{0pt}{1.5em}{#1}{#2}% ← leader dots
  \vspace{4pt}% ← space after chapter entry
}

% Section dividers
\def\l@sectiondivider#1#2{%
  \vspace{20pt}% before
  \ifnum \c@tocdepth >\m@ne
    \addpenalty{-\@highpenalty}%
    \vskip 1em \@plus\p@
    \setlength\@tempdima{1.5em}%
    \begingroup
      \parindent \z@ \rightskip \@pnumwidth
      \parfillskip -\@pnumwidth
      \leavevmode \fontsize{12pt}{14pt}\selectfont\bfseries\color{sectioncolor}
      \advance\leftskip\@tempdima
      \hskip -\leftskip
      #1\nobreak\hfil \nobreak\hb@xt@\@pnumwidth{\hss #2}\par
    \endgroup
  \fi
  \vspace{4pt}% after
}

% Generic dividers
\def\l@genericdivider#1#2{%
  \vspace{6pt}% before
  \ifnum \c@tocdepth >\m@ne
    \addpenalty{-\@highpenalty}%
    \vskip 0.5em \@plus\p@
    \setlength\@tempdima{1.5em}%
    \begingroup
      \parindent \z@ \rightskip \@pnumwidth
      \parfillskip -\@pnumwidth
      \leavevmode \fontsize{11pt}{13pt}\selectfont\bfseries\color{accentcolor}
      \advance\leftskip\@tempdima
      \hskip -\leftskip
      #1\nobreak\hfil \nobreak\hb@xt@\@pnumwidth{\hss #2}\par
    \endgroup
  \fi
  \vspace{2pt}% after
}
\makeatother

% Default fancy page style (Regular pages)
\fancypagestyle{fancy}{%
    \fancyhf{}
    \fancyhead[L]{\footnotesize \textit{Visual Layer for Camtek}}
    % Header Right displays current chapter/appendix name
    \fancyhead[R]{\footnotesize \leftmark}

    % Footer: White text, no line, PUSHED DOWN using raisebox
    % Note: These settings must match the 'plain' style below exactly
    \fancyfoot[L]{\raisebox{-9pt}{\footnotesize \raisebox{-6pt}{
\includegraphics[height=0.6cm]{images/logo.png}}}}
    \fancyfoot[C]{\raisebox{-10pt}{\footnotesize \color{white} \textit{Last updated:} \today}}
    \fancyfoot[R]{\raisebox{-10pt}{\footnotesize \color{white} \thepage}}

    \renewcommand{\headrulewidth}{0.4pt}
    \renewcommand{\footrulewidth}{0pt} % Removes the line above the footer
}

% Plain style (First page of chapters) - Updated to match 'fancy' footer
\fancypagestyle{plain}{%
    \fancyhf{}
    % No header content or rule for chapter start pages
    \renewcommand{\headrulewidth}{0pt}

    % Footer: Exact copy of 'fancy' footer to ensure consistency
    \fancyfoot[L]{\raisebox{-9pt}{\footnotesize \raisebox{-6pt}{
\includegraphics[height=0.6cm]{images/logo.png}}}}
    \fancyfoot[C]{\raisebox{-10pt}{\footnotesize \color{white} \textit{Last updated:} \today}}
    \fancyfoot[R]{\raisebox{-10pt}{\footnotesize \color{white} \thepage}}

    \renewcommand{\footrulewidth}{0pt} % Removes the line above the footer
}

% Empty style for section divider pages (no header, no footer)
\fancypagestyle{empty}{%
  \fancyhf{}
  \renewcommand{\headrulewidth}{0pt}
  \renewcommand{\footrulewidth}{0pt}
}

% TOC style (no header, keep footer)
\fancypagestyle{tocstyle}{%
    \fancyhf{}
    % No header content or rule for chapter start pages
    \renewcommand{\headrulewidth}{0pt}

    % Footer: Exact copy of 'fancy' footer to ensure consistency
    \fancyfoot[L]{\raisebox{-9pt}{\footnotesize \raisebox{-6pt}{
\includegraphics[height=0.6cm]{images/logo.png}}}}
    \fancyfoot[C]{\raisebox{-10pt}{\footnotesize \color{white} \textit{Last updated:} \today}}
    \fancyfoot[R]{\raisebox{-10pt}{\footnotesize \color{white} \thepage}}

    \renewcommand{\footrulewidth}{0pt} % Removes the line above the footer
}

% Customize TOC section entries
\usepackage{tocloft}

% TOC Entry Hierarchy:
% 1. Section Dividers (numbered) - "Section I: Getting Started", "Section II: Main Workflow"
%    Style: Large (12pt), bold, sectioncolor, extra spacing
% 2. Appendix Divider (unnumbered) - "Appendices"
%    Style: Large (12pt), bold, sectioncolor, extra spacing (same as section dividers)
% 3. Generic Dividers (unnumbered) - "Step 1: Create & Prepare Dataset"
%    Style: Medium (11pt), bold, emphasiscolor, moderate spacing
% 4. Chapters (regular numbered) - Default report class styling
%    Style: Normal chapter formatting with numbers

% Style for numbered section dividers (Section I: Getting Started, etc.)
\makeatletter
\def\l@sectiondivider#1#2{%
  \ifnum \c@tocdepth >\m@ne
    \addpenalty{-\@highpenalty}%
    \vskip 2em \@plus\p@  % Extra space before
    \setlength\@tempdima{1.5em}%
    \begingroup
      \parindent \z@ \rightskip \@pnumwidth
      \parfillskip -\@pnumwidth
      \leavevmode \bfseries
      \fontsize{12pt}{14pt}\selectfont
      \color{sectioncolor}%
      \advance\leftskip\@tempdima
      \hskip -\leftskip
      #1\nobreak\hfil \nobreak\hb@xt@\@pnumwidth{\hss #2}\par
      \penalty\@highpenalty
    \endgroup
    \vskip 0.5em \@plus\p@  % Space after
  \fi
}

% Style for appendix divider (same as section dividers but unnumbered)
\def\l@appendixdivider#1#2{%
  \ifnum \c@tocdepth >\m@ne
    \addpenalty{-\@highpenalty}%
    \vskip 2em \@plus\p@  % Extra space before
    \setlength\@tempdima{1.5em}%
    \begingroup
      \parindent \z@ \rightskip \@pnumwidth
      \parfillskip -\@pnumwidth
      \leavevmode \bfseries
      \fontsize{12pt}{14pt}\selectfont
      \color{sectioncolor}%
      \advance\leftskip\@tempdima
      \hskip -\leftskip
      #1\nobreak\hfil \nobreak\hb@xt@\@pnumwidth{\hss #2}\par
      \penalty\@highpenalty
    \endgroup
    \vskip 0.5em \@plus\p@  % Space after
  \fi
}

% Style for unnumbered generic dividers (Step 1, Step 2, etc.)
\def\l@genericdivider#1#2{%
  \ifnum \c@tocdepth >\m@ne
    \addpenalty{-\@highpenalty}%
    \vskip 1.0em \@plus\p@  % Medium space before
    \setlength\@tempdima{1.5em}%
    \begingroup
      \parindent \z@ \rightskip \@pnumwidth
      \parfillskip -\@pnumwidth
      \leavevmode \bfseries
      \fontsize{11pt}{13pt}\selectfont
      \color{emphasiscolor}%
      \advance\leftskip\@tempdima
      \hskip -\leftskip
      #1\nobreak\hfil \nobreak\hb@xt@\@pnumwidth{\hss #2}\par
      \penalty\@highpenalty
    \endgroup
    \vskip 0.3em \@plus\p@  % Space after
  \fi
}
\makeatother

% Set default page style
\pagestyle{fancy}
\setlength{\footskip}{64pt}

% Make chapter pages use plain style (no header)
\patchcmd{\chapter}{\thispagestyle{plain}}{\thispagestyle{plain}}{}{}

% Apply tocstyle to table of contents
\usepackage{etoolbox}
\preto\tableofcontents{%
  \clearpage
  \thispagestyle{tocstyle}%
  \pagestyle{tocstyle}%
  \addtocontents{toc}{\protect\widowpenalty=10000}%
  \addtocontents{toc}{\protect\clubpenalty=10000}%
  \addtocontents{toc}{\protect\interlinepenalty=10000}%
}
\appto\tableofcontents{\pagestyle{fancy}}

% Prevent page breaks between TOC entries and their sub-entries
\makeatletter
\renewcommand{\@starttoc}[1]{%
  \begingroup
    \@input{\jobname.#1}%
    \if@filesw
      \expandafter\newwrite\csname tf@#1\endcsname
      \immediate\openout \csname tf@#1\endcsname \jobname.#1\relax
    \fi
    \@nobreakfalse
  \endgroup
}
\makeatother

% Table styling
% NOTE: array, longtable are loaded by Pandoc, but we need to ensure booktabs is loaded
% before we redefine its commands later
\usepackage{array}
\usepackage{booktabs}
\usepackage{longtable}
\usepackage{tabularx}
\usepackage{colortbl}

% Create new column type M for middle-aligned columns (instead of redefining p)
% This avoids conflicts with standard p column type
\newcolumntype{M}[1]{m{#1}}

% Define flexible column types for tabularx (auto-expanding to fill width)
\newcolumntype{Y}{>{\centering\arraybackslash}X}  % Centered, auto-width
\newcolumntype{Z}{>{\raggedright\arraybackslash}X} % Left-aligned, auto-width


% Apply font sizing to tabularx environment
\AtBeginEnvironment{tabularx}{%
  \fontsize{9pt}{10pt}\selectfont
}


% Note: toprule, midrule, bottomrule are redefined later with bold formatting

% Define admonition colors
\definecolor{notecolor}{RGB}{214, 214, 214} % #D6D6D6
\definecolor{importantcolor}{RGB}{255, 154, 3} % #FF9A03
\definecolor{warningcolor}{RGB}{255, 105, 94} % #FF695E
\definecolor{tipcolor}{RGB}{51, 108, 155} % #336C9B
\definecolor{criticalcolor}{RGB}{255, 105, 94} % #FF695E
\definecolor{plaintable}{RGB}{244, 242, 248}
\definecolor{darkfont}{RGB}{21, 6, 43}
\definecolor{imageborder}{HTML}{1D102D}



% Process openings
\newcommand{\processstart}{%
  \needspace{4\baselineskip}%  Ensure at least 4 lines of space before starting
  \vspace{6pt}%
  \par\noindent%
  \begin{minipage}{\textwidth}%  Use minipage instead of table to prevent floating
  \nopagebreak[4]%  Strongly discourage page break after this
  \rowcolors{1}{}{}%
  \arrayrulecolor{sectioncolor}%
  \begin{tabular}{>{\color{sectioncolor}\bfseries\setlength{\parindent}{-7pt}\setlength{\leftskip}{2pt}\setlength{\rightskip}{8pt}\fontsize{13pt}{14.5pt}\selectfont}p{0.7\textwidth}}%
  \nopagebreak[4]%  Prevent break between tabular start and content
}
\newcommand{\processend}{%
  \nopagebreak[4]%  Prevent break before end
  \\ \hline
  \end{tabular}%
  \arrayrulecolor{tableborder}%
  \end{minipage}%
  \par%
  \vspace{5pt}%
}





% Admonition table styling
% Note admonitions
\newcommand{\notestart}{%
  \needspace{3\baselineskip}%
  \vspace{6pt}%
  \nopagebreak[4]%
  \begin{table}[!ht]
  \centering
  \rowcolors{1}{}{}
  \arrayrulecolor{notecolor}
  \begin{tabular}{|>{\centering\arraybackslash}m{1cm}|>{\columncolor{notecolor!20}\setlength{\parindent}{0pt}\setlength{\leftskip}{6pt}\setlength{\rightskip}{8pt}\fontsize{8pt}{9.7pt}\selectfont}p{0.63\textwidth}|}
  \hline
  \cellcolor{notecolor!30}\raisebox{-0.55\height}{
\includegraphics[height=0.5cm]{images/information-note.png}} &
}
\newcommand{\noteend}{%
  \\ \hline
  \end{tabular}

  \end{table}
  \arrayrulecolor{tableborder}
  \vspace{3pt}
}

% Important admonitions
\newcommand{\importantstart}{%
  \needspace{3\baselineskip}%
  \vspace{6pt}%
  \nopagebreak[4]%
  \begin{table}[!ht]
  \centering
  \rowcolors{1}{}{}
  \arrayrulecolor{importantcolor}
  \begin{tabular}{|>{\centering\arraybackslash}m{1cm}|>{\columncolor{importantcolor!15}\setlength{\parindent}{0pt}\setlength{\leftskip}{6pt}\setlength{\rightskip}{8pt}\fontsize{8pt}{9.7pt}\selectfont}p{0.63\textwidth}|}
  \hline
  \cellcolor{importantcolor!17}\raisebox{-0.55\height}{
\includegraphics[height=0.5cm]{images/important.png}} &
}

\newcommand{\importantend}{%
  \\ \hline
  \end{tabular}
  \end{table}
  \arrayrulecolor{tableborder}
  \vspace{6pt}
}

% Warning admonitions
\newcommand{\warningstart}{%
  \needspace{3\baselineskip}%
  \vspace{6pt}%
  \nopagebreak[4]%
  \begin{table}[!ht]
  \centering
  \rowcolors{1}{}{}
  \arrayrulecolor{warningcolor}
  \begin{tabular}{|>{\centering\arraybackslash}m{1cm}|>{\columncolor{warningcolor!17}\setlength{\parindent}{0pt}\setlength{\leftskip}{6pt}\setlength{\rightskip}{8pt}\fontsize{8pt}{9.7pt}\selectfont}p{0.63\textwidth}|}
  \hline
  \cellcolor{warningcolor!17}\raisebox{-0.55\height}{
\includegraphics[height=0.5cm]{images/warning-critical.png}} &
}
\newcommand{\warningend}{%
  \\ \hline
  \end{tabular}
  \end{table}
  \arrayrulecolor{tableborder}
  \vspace{6pt}
}

% Tip admonitions
\newcommand{\tipstart}{%
  \needspace{3\baselineskip}%
  \vspace{6pt}%
  \nopagebreak[4]%
  \begin{table}[!ht]
  \centering
  \rowcolors{1}{}{}
  \arrayrulecolor{tipcolor}
  \begin{tabular}{|>{\centering\arraybackslash}m{1cm}|>{\columncolor{tipcolor!15}\setlength{\parindent}{0pt}\setlength{\leftskip}{6pt}\setlength{\rightskip}{8pt}\fontsize{8pt}{9.7pt}\selectfont}p{0.63\textwidth}|}
  \hline
  \cellcolor{tipcolor!17}\raisebox{-0.55\height}{
\includegraphics[height=0.5cm]{images/tip.png}} &
}
\newcommand{\tipend}{%
  \\ \hline
  \end{tabular}
  \end{table}
  \arrayrulecolor{tableborder}
  \vspace{6pt}
}

% Critical admonitions
\newcommand{\criticalstart}{%
  \needspace{3\baselineskip}%
  \vspace{6pt}%
  \nopagebreak[4]%
  \begin{table}[!ht]
  \centering
  \rowcolors{1}{}{}
  \arrayrulecolor{criticalcolor}
  \begin{tabular}{|>{\centering\arraybackslash}m{1cm}|>{\columncolor{criticalcolor!15}\setlength{\parindent}{0pt}\setlength{\leftskip}{6pt}\setlength{\rightskip}{8pt}\fontsize{8pt}{9.7pt}\selectfont}p{0.63\textwidth}|}
  \hline
  \cellcolor{criticalcolor!17}\raisebox{-0.55\height}{
\includegraphics[height=0.5cm]{images/warning-critical.png}} &
}
\newcommand{\criticalend}{%
  \\ \hline
  \end{tabular}
  \end{table}
  \arrayrulecolor{tableborder}
  \vspace{4pt}
}

% Enhanced table styling with header detection
% (oldtabular/oldendtabular no longer needed - using AtBeginEnvironment instead)



% Define light grey color for table borders (for regular tables only, not admonitions)
\definecolor{tableborder}{RGB}{81, 67, 132}
\definecolor{tableprimary}{RGB}{233, 228, 241}
\definecolor{tablesecondary}{RGB}{244, 241, 248}

% Plain table environment with custom styling (for markdown tables)
% Usage in markdown: tables are automatically wrapped
% Styling: light grey borders, bold headers, clean appearance
\newenvironment{plaintable}[1]{%
  \begin{table}[h]
  \centering
  \setlength{\arrayrulewidth}{1pt}
  \rowcolors{1}{tableprimary}{tablesecondary}
  \begin{tabular}{#1}
    
}{%
  \end{tabular}
  \end{table}
}

% Apply font styling to tabular environment
\AtBeginEnvironment{tabular}{%
  \fontsize{10pt}{12pt}\selectfont
}

% Code block styling with enhanced JSON highlighting
\usepackage{listings}

\definecolor{codebg}{RGB}{81, 67, 132}
\definecolor{backcolour}{rgb}{0.96,0.95,0.97}
\definecolor{jsonkey}{RGB}{88, 110, 117}
\definecolor{jsonstring}{RGB}{42, 161, 152}
\definecolor{jsonnumber}{RGB}{211, 54, 130}

% Define JSON language for better syntax highlighting
\lstdefinelanguage{json}{
    string=[s]{"}{"},
    stringstyle=\color{jsonstring},
    comment=[l]{:},
    commentstyle=\color{jsonkey},
    morestring=[b]',
    literate=
     *{0}{{{\color{jsonnumber}0}}}{1}
      {1}{{{\color{jsonnumber}1}}}{1}
      {2}{{{\color{jsonnumber}2}}}{1}
      {3}{{{\color{jsonnumber}3}}}{1}
      {4}{{{\color{jsonnumber}4}}}{1}
      {5}{{{\color{jsonnumber}5}}}{1}
      {6}{{{\color{jsonnumber}6}}}{1}
      {7}{{{\color{jsonnumber}7}}}{1}
      {8}{{{\color{jsonnumber}8}}}{1}
      {9}{{{\color{jsonnumber}9}}}{1}
}

\lstdefinestyle{codebg}{
    backgroundcolor=\color{snippets},
    basicstyle=\ttfamily\normalsize,
    breakatwhitespace=true,
    breaklines=true,
    captionpos=b,
    keepspaces=true,
    numbersep=5pt,
    showspaces=false,
    showstringspaces=false,
    showtabs=false,
    tabsize=2,
    language=json
}

\lstset{style=codebg}

% Regular table header styling with 900 weight
\newcommand{\tableheader}[1]{\fontsize{9pt}{10pt}\selectfont}

% Apply plain table styling to longtable (for Pandoc markdown tables)
\AtBeginEnvironment{longtable}{%
  \setlength{\arrayrulewidth}{1pt}
  % First row (header) gets tableprimary, then alternate tablesecondary and white
  \rowcolors{1}{tableprimary}{tablesecondary}
  % Set font size to 10pt for all table content
  \fontsize{9pt}{10pt}\selectfont
  % Make longtable span full text width
  \setlength{\LTleft}{0pt}
  \setlength{\LTright}{0pt}
}

% Redefine booktabs rules to be simple hlines
\renewcommand{\toprule}{\hline}
\renewcommand{\midrule}{\hline}
\renewcommand{\bottomrule}{\hline}

% Set borders for all regular tables
\setlength{\arrayrulewidth}{0.7pt}   % sets all \hline and vertical rules to 0.5pt
\setlength{\tabcolsep}{4pt}           % optional: spacing inside cells
\renewcommand{\arraystretch}{1.3}     % maintain row padding
\arrayrulecolor{tableborder}     % sets color for all borders

\usepackage{bookmark}
\IfFileExists{xurl.sty}{\usepackage{xurl}}{} % add URL line breaks if available
\urlstyle{same}
\hypersetup{
  colorlinks=true,
  linkcolor={linkcolor},
  filecolor={Maroon},
  citecolor={blue},
  urlcolor={blue},
  pdfcreator={LaTeX via pandoc}}

\author{}
\date{}

\begin{document}

{
\hypersetup{linkcolor=}
\setcounter{tocdepth}{2}
\tableofcontents
}
\setstretch{1.2}
\chapter{About This Documentation}\label{about-this-documentation}

This documentation is specifically designed for the Camtek Visual Layer
implementation. Camtek uses Visual Layer for manufacturing quality
control, where the system processes output to rapidly train defect
detection models. This implementation is deployed on-premises, allowing
direct access to machine output folders rather than traditional cloud
uploads.

\sectiondivider{An Overview of Visual Layer}

\chapter{\texorpdfstring{Integrating \textbf{Visual Layer} in
Manufacturing
Workflows}{Integrating Visual Layer in Manufacturing Workflows}}\label{integrating-visual-layer-in-manufacturing-workflows}

\textbf{Visual Layer} is a unified platform, designed to help teams
manage, curate, and analyze visual data at scale, streamlining the
entire data lifecycle. The platform provides centralized tools for
exploring datasets, identifying patterns, maintaining data quality, and
collaborating across teams.

\tipstart

\textbf{New to Visual Layer terminology?} A comprehensive glossary with
detailed definitions, examples, and context for all concepts can be
found in the \hyperref[key-concepts]{Key Concepts} section of the
Appendices. \tipend

Key capabilities for Camtek users include:

\begin{itemize}
\tightlist
\item
  Automated human-in-the-loop labeling
\item
  Performance analysis and tracking
\item
  Integration with the Camtek training infrastructure for real-time
  model deployment
\end{itemize}

The rest of this section describes:

\begin{itemize}
\tightlist
\item
  \hyperref[key-concepts]{Key concepts}
\item
  \hyperref[browser-requirements]{Prerequisites for exploring datasets
  with Visual Layer}
\end{itemize}

\section{Key Concepts}\label{key-concepts}

Before exploring \textbf{Visual Layer}, familiarize yourself with these
core concepts:

\begin{longtable}[]{@{}
  >{\raggedright\arraybackslash}p{(\linewidth - 2\tabcolsep) * \real{0.5000}}
  >{\raggedright\arraybackslash}p{(\linewidth - 2\tabcolsep) * \real{0.5000}}@{}}
\toprule\noalign{}
\begin{minipage}[b]{\linewidth}\raggedright
\textbf{Concept}
\end{minipage} & \begin{minipage}[b]{\linewidth}\raggedright
\textbf{Description}
\end{minipage} \\
\midrule\noalign{}
\endhead
\bottomrule\noalign{}
\endlastfoot
\textbf{Class} & Each distinct type of label in your classification
system. You need at least two classes, each with enough examples. \\
\textbf{Class Taxonomy} & The complete structure of all possible classes
defined in classes.json from your AOI machine output. Cannot be changed
after dataset creation. \\
\textbf{Cluster} & A group of visually similar images automatically
organized by \textbf{Visual Layer} \\
\textbf{Confidence Score} & A numerical measure (0-100\%) indicating how
certain the system is about an automatically assigned label. \\
\textbf{Dataset} & A collection of images and associated metadata
organized for analysis and labeling \\
\textbf{Embedding} & The technology that enables \textbf{Visual Layer}
to find visually similar images and automatically label thousands from
just a few examples. \\
\textbf{Ground Truth} & Verified, accurate labels serving as the
authoritative reference for training and evaluating machine learning
models. \\
\textbf{Label} & The category assigned to a single image---the answer to
``what type is this?'' \\
\textbf{Label Propagation} & Semi-supervised learning combining
automated pattern recognition with human oversight to efficiently label
large datasets. \\
\textbf{Seed} & Initial labeled examples provided for each class. These
training examples teach the system what each class looks like. \\
\end{longtable}

\tipstart

\textbf{Need more detail?} For comprehensive definitions with examples
and usage context, see the \hyperref[key-concepts]{Key Concepts} section
in the Appendices. \tipend

\section{Prerequisites}\label{browser-requirements}

Chrome browser required.

\sectiondivider{The Workflow}

\chapter{The Workflow: Getting
Started}\label{the-workflow-getting-started}

This section guides you through the complete \textbf{Visual Layer}
workflow for manufacturing defect detection, from dataset creation
through model training.

The following diagram outlines the recommended sequence for using
\textbf{Visual Layer} within a standard defect-detection workflow.

\begin{figure}[H]
        \includegraphics[width=.8\textwidth]{images/vl-manufacturing-workflow-overview-diagram.png}
        \caption{The Primary Workflow}
        \label{fig:wrapfig}
\end{figure}

The key steps of this workflow are as follows:

\renewcommand{\arraystretch}{1.5}
\rowcolors{1}{tableprimary}{tablesecondary}

\begin{tabularx}{\textwidth}{|>{\centering\arraybackslash}m{4em}|p{0.35\textwidth}|X|}
\hline
\textbf{Step} & \textbf{Phase} & \textbf{Description} \\
\hline

\raisebox{-7pt}{
\includegraphics[width=5.5mm]{images/1light.png}} & \textbf{\hyperref[how-dataset-creation-works]{Create \& Prepare Dataset}} & Ingest your production imagery and metadata from manufacturing equipment into the \textbf{Visual Layer} platform, then explore it to prepare for label propagation. \\
\hline

\raisebox{-7pt}{
\includegraphics[width=5.5mm]{images/2light.png}} & \textbf{\hyperref[how-label-propagation-works]{Label Propagation}} & Add representative samples, automatically propagate labels to similar images, and review and validate labels to ensure consistency and accuracy across your dataset. \\
\hline

\raisebox{-7pt}{
\includegraphics[width=5.5mm]{images/3light.png}} & \textbf{\hyperref[train-models-how-it-works]{Train Models Downstream}} & Send your finalized, validated dataset to your model-training pipeline for deployment. \\
\hline
\end{tabularx}

\genericdivider{Step 1: Create \& Prepare Dataset}

\chapter{Create \& Prepare: How It
Works}\label{how-dataset-creation-works}

Creating a dataset is the first main step in the \textbf{Visual Layer}
workflow for curating your data and labeling images.

Camtek automated optical inspection (AOI) scans produce a standardized
directory containing inspection results, production metadata,
class/label definitions, and the associated images organized into
subfolders. \textbf{Visual Layer} processes the output from Camtek AOI
machines to rapidly train defect detection models.

\notestart

Multiple folders must share identical Job/Setup/Recipe to be combined.
This structure is required in order to ensure that in any given dataset,
only outputs from the same production job are included. \noteend

You can quickly create and prepare a dataset for label propagation as
outlined in the following diagram:

\begin{figure}[H]
\centering
\includegraphics[width=0.8\textwidth]{images/vl-create-prepare-flow-diagram.png}
\caption{Dataset Creation \& Preparation Workflow}
\end{figure}

\renewcommand{\arraystretch}{1.5}
\rowcolors{1}{tableprimary}{tablesecondary}
\begin{tabularx}{\textwidth}{|>{\centering\arraybackslash}m{4em}|p{0.35\textwidth}|X|}
\hline
\textbf{Step} & \textbf{Phase} & \textbf{Description} \\
\hline
\raisebox{-7pt}{
\includegraphics[width=5.5mm]{images/1light.png}} & \textbf{Upload \& Select Folders} & Browse and upload AOI output directories that you want to include in the dataset. Upload as many as you'd like. For each folder, \textbf{Visual Layer} then automatically validates required files.\\
\hline
\raisebox{-7pt}{
\includegraphics[width=5.5mm]{images/2light.png}} & \textbf{Build \& Create} & From the list of uploaded folders, checkmark all the folders to be included in the dataset. \textbf{Visual Layer} then creates the dataset, indexing and clustering images, and loading your class taxonomy based on the information that was included in each uploaded folder. Classes become permanent after creation and cannot be modified.  \\
\hline
\raisebox{-7pt}{
\includegraphics[width=5.5mm]{images/3light.png}} & \textbf{Explore Dataset} & When the \hyperref[dataset-creation-stages]{dataset is ready}, explore the dataset to identify representative examples for each defect class.  \\
\hline
\end{tabularx}

\chapter{Create \& Prepare: Upload \&
Build}\label{creating-datasets-from-machine-output}

This section guides you through creating datasets in \textbf{Visual
Layer} from one or more AOI machine output directories derived from the
same production job.

You'll:

\begin{itemize}
\tightlist
\item
  choose and upload output folders from your machines for Visual Layer
  validation
\item
  select the folders relevant for one specific dataset
\item
  confirm and wait for Visual Layer to generate the dataset
\end{itemize}

\processstart

\textbf{To create a dataset} \processend

\begin{enumerate}
\def\labelenumi{\arabic{enumi}.}
\item
  From \textbf{Dataset Inventory} (the ``Home'' page) click \textbf{New
  Dataset.}

  The \textbf{Data} area loads.
\item
  From the \textbf{Enter Folder Paths} area, enter only the name of the
  relevant folder you want to upload.
\end{enumerate}

\notestart The base path to your machine output is configured during
installation. Enter only the folder name. \noteend

\begin{figure}[H]
   \centering
   \includegraphics[width=0.7\textwidth]{images/vl-create-dataset-onprem.png}
   \caption{Create a New Dataset}
   \end{figure}

\begin{enumerate}
\def\labelenumi{\arabic{enumi}.}
\setcounter{enumi}{2}
\item
  You can upload additional folders at the same time, or one at a time.
  To add an additional folder, click \textbf{+}.

  An additional field appears for another folder name. Add as many as
  you need.
\item
  Click \textbf{Add} to add the folder/s to the list.

  \textbf{Visual Layer} extracts metadata and validates folder structure
  and contents. When complete, the folders appear in the
  \textbf{Selected Folders} table.

  \begin{figure}[H]
  \centering
  \includegraphics[width=0.8\textwidth]{images/vl-create-dataset-selectfolders.png}
  \caption{Selected Folders table}
  \end{figure}

  The table displays the following information for each folder:

  \begin{center}
  \begin{tabular}{|l|l|}
  \hline
  \textbf{Column} & \textbf{Description} \\
  \hline
  \textbf{Job Name} & The production job identifier from the folder metadata \\
  \hline
  \textbf{Setup Name} & The setup configuration identifier from the folder metadata \\
  \hline
  \textbf{Recipe List} & The recipe name(s) used during inspection \\
  \hline
  \textbf{Image Count} & Total number of images contained in the folder \\
  \hline
  \end{tabular}
  \end{center}

  \importantstart

  All folders in a dataset must share identical Job Name, Setup Name,
  and Recipe List. \textbf{Visual Layer} validates compatibility across
  all folders to ensure they originate from the same production run.
  Folders with different metadata are automatically disabled and grayed
  out. \importantend
\item
  If you uploaded multiple folders and need to identify which are
  compatible, use the table filters:

  \begin{itemize}
  \item
    \textbf{Filter by production metadata} - Use the Job Name, Setup
    Name, and Recipe Name filters to find folders from the same
    production run:

    \includegraphics[width=0.7\linewidth,height=\textheight,keepaspectratio]{images/vl-create-dataset-selectfolders-filters.png}
  \item
    \textbf{Check compatibility} - Hover over the info icon to view
    details about any incompatible folders:

    \includegraphics[width=0.7\linewidth,height=\textheight,keepaspectratio]{images/vl-create-dataset-selectfolders-error.png}
  \end{itemize}
\item
  Select the folders to include in the dataset.

  When you select any folder, all folders that are not part of the same
  job, recipe and setup are filtered out of the list automatically. To
  see those folders again, uncheck any checked folders.
\item
  Click \textbf{Continue}.

  \textbf{Visual Layer} loads a preview of the intended dataset (a
  combination of all of the folders you selected) for you to review and
  approve.

  \begin{figure}[H]
  \centering
  \includegraphics[width=0.6\textwidth]{images/vl-create-dataset-preview.png}
  \caption{Preview Dataset}
  \end{figure}
\item
  Click \textbf{Create Dataset}.

  \textbf{Visual Layer} begins processing the dataset.

  \begin{figure}[H]
  \centering
  \includegraphics[width=0.6\textwidth]{images/vl-create-dataset-progress.png}
  \caption{Creation Progress steps}
  \end{figure}

  Processing time varies with dataset size, typically 5-15 minutes for
  standard production runs. When complete, you'll receive a
  notification.
\end{enumerate}

\chapter{Create \& Prepare: Exploration}\label{exploration}

Once your dataset is ready, spend time exploring it before beginning
label propagation to identify representative examples for each defect
class. This exploration phase is critical for selecting high-quality
seed examples that will train the system effectively.

\tipstart

Spend sufficient time exploring before starting label propagation. The
quality of your seed examples directly impacts automatic labeling
accuracy. Look for clear, unambiguous examples that represent the full
visual range of each defect class. \tipend

\processstart

\textbf{To explore and prepare for label propagation} \processend

\begin{enumerate}
\def\labelenumi{\arabic{enumi}.}
\item
  From the \textbf{Dataset Inventory}, search or scroll to find your
  dataset and click its card to open it.

  The dataset loads in cluster view, displaying groups of visually
  similar images.

  \begin{figure}
  \centering
  \includegraphics[width=0.9\linewidth,height=\textheight,keepaspectratio,alt={Dataset Cluster View}]{images/vl-overview-cluster-view.png}
  \caption{Dataset Cluster View}
  \end{figure}

  \textbf{Visual Layer} automatically organizes your images into
  clusters during dataset creation. Each cluster groups visually similar
  images together and displays as a preview card showing:

  \begin{itemize}
  \tightlist
  \item
    \textbf{Number of items} - Total images in the cluster matching your
    current filters
  \item
    \textbf{Representative thumbnails} - Key example images that
    summarize the cluster's visual patterns
  \item
    \textbf{Prominent labels} - Most frequent labels within the cluster
    (if labels exist)
  \end{itemize}

  \begin{figure}
  \centering
  \includegraphics[width=0.93\linewidth,height=\textheight,keepaspectratio,alt={Cluster preview cards}]{images/vl-clusters-image-view.png}
  \caption{Cluster preview cards}
  \end{figure}

  Clusters help you quickly identify patterns and find similar defects
  without examining every image individually.
\item
  Look for clusters that clearly represent distinct defect types. Click
  any cluster to examine its images more closely and understand what
  visual patterns define each class.
\item
  Click
  
\includegraphics[width=0.05\linewidth,height=\textheight,keepaspectratio]{images/vl-filterbutton.png}
  to open the \textbf{Filter Menu} and narrow down to specific subsets:

  \begin{itemize}
  \tightlist
  \item
    Filter by \textbf{Folders} to examine specific production runs
  \item
    Use \textbf{Quality Issues} to exclude poor-quality images
  \item
    Apply \textbf{Duplicates} filter to avoid redundant examples
  \end{itemize}

  \begin{figure}
  \centering
  \includegraphics[width=0.4\linewidth,height=\textheight,keepaspectratio,alt={All Filters}]{images/vl-camtek-allfilters.png}
  \caption{All Filters}
  \end{figure}
\item
  When you find a good example of a defect class, hover over the image
  and click
  
\includegraphics[width=0.05\linewidth,height=\textheight,keepaspectratio]{images/vl-visual-similarity-icon.png}
  to find similar images. This helps you locate all potential candidates
  for that class.
\item
  Based on your exploration, take note which defect classes are
  well-represented in your data and where you can find at least 5 clear,
  diverse examples for each class.
\end{enumerate}

\processstart

\textbf{Next Steps} \processend

Now that you've explored your dataset and identified representative
examples for each defect class, you're ready to begin
\hyperref[how-label-propagation-works]{label propagation}.

\genericdivider{Step 2: Label Propagation}

\chapter{Label Propagation: How It
Works}\label{how-label-propagation-works}

Label propagation is a semi-automated labeling process that transforms
visual data annotation from a one-time task into a continuous
improvement cycle. \textbf{Visual Layer} first receives instructions and
human oversight to teach the system which labels to apply and for what
kinds of images, combining automated pattern recognition with human
oversight to create a feedback loop where each iteration improves
labeling accuracy.

This feedback loop gives the system the guidance necessary in order to
guarantee near-perfect automatic labeling thereafter.

\notestart

Each dataset maintains its own independent model, preventing
cross-contamination between production runs. \noteend

\begin{figure}[H]
\setlength{\leftskip}{1pt}
\includegraphics[width=.9\textwidth]{images/vl-labelpropagation-workflow-diagram.png}
\caption{Label Propagation Workflow}
\end{figure}

\newpage

\textbf{The label propagation process includes these substeps:}

\renewcommand{\arraystretch}{1.5}
\rowcolors{1}{tableprimary}{tablesecondary}

\begin{tabularx}{\textwidth}{|>{\centering\arraybackslash}m{4em}|p{0.35\textwidth}|X|}
\hline
\textbf{Step} & \textbf{Phase} & \textbf{Description} \\
\hline

\raisebox{-7pt}{
\includegraphics[width=5.5mm]{images/1light.png}}
  & \textbf{Seed Creation}
  & Provide ≥5 examples per class to teach the system what each class looks like. \\
\hline

\raisebox{-7pt}{
\includegraphics[width=5.5mm]{images/2light.png}}
  & \textbf{Iterative Review Cycle}
  & Work with \textbf{Visual Layer} to refine labels:
    \begin{itemize}
        \item The system automatically labels high-confidence images
        \item Uncertain images are sent to you in small review batches
        \item You review and provide feedback on these batches
        \item The system learns and improves with each review round
    \end{itemize} \\
\hline

\raisebox{-7pt}{
\includegraphics[width=5.5mm]{images/3light.png}}
  & \textbf{Automatic Completion}
  & When model confidence reaches the threshold, label propagation automatically completes and applies final labels to the entire dataset. \\
\hline
\end{tabularx}

\chapter{Label Propagation: Seed Creation}\label{seed-creation}

To get started with label propagation you will provide \textbf{Visual
Layer} with seed examples that train our system to automatically label
your remaining data. Label propagation can run on either the entire
dataset or a subset, for example newly added data only.

All labels during label propagation are stored in a temporary table and
can be edited. When you click \textbf{Finish} or the process completes
naturally, the labels are written permanently to your dataset. After
finalization, labels cannot be edited without running a new label
propagation iteration.

\importantstart

\textbf{Visual Layer} assigns a single label to each image during label
propagation. Multi-label classification is not supported. \newline

Starting label propagation over from scratch \textbf{\textbf{replaces}}
all previous label propagation results---existing labels are overwritten
and cannot be recovered. \importantend

\processstart

\textbf{To maximize success with label propagation} \processend

\begin{longtable}[]{@{}
  >{\raggedright\arraybackslash}p{(\linewidth - 2\tabcolsep) * \real{0.5000}}
  >{\raggedright\arraybackslash}p{(\linewidth - 2\tabcolsep) * \real{0.5000}}@{}}
\toprule\noalign{}
\begin{minipage}[b]{\linewidth}\raggedright
\textbf{Best Practice}
\end{minipage} & \begin{minipage}[b]{\linewidth}\raggedright
\textbf{Description}
\end{minipage} \\
\midrule\noalign{}
\endhead
\bottomrule\noalign{}
\endlastfoot
\textbf{Plan Your Classes} & Select ≥2 classes. Include every class that
may appear---classes cannot be added later. \\
\textbf{Provide ≥5 Examples Per Class} & Minimum required to proceed.
More diverse examples = better accuracy. \\
\textbf{Choose Representative Examples} & Select clear, typical
instances showing the full visual range of each class. Avoid edge cases
or ambiguous images. \\
\textbf{Triple-Check Before Starting} & Seed examples cannot be edited
once propagation begins. Verify all assignments are correct. \\
\end{longtable}

\processstart

\textbf{To prepare your dataset} \processend

\notestart

During seed creation, only class assignment is available and other
exploration features are disabled. \noteend

\begin{enumerate}
\def\labelenumi{\arabic{enumi}.}
\item
  Navigate to the relevant dataset in \textbf{Visual Layer}.
\item
  From the panel on the right side, click
  \includegraphics[width=0.05\linewidth,height=\textheight,keepaspectratio]{images/vl-labelpropagation-tab-icon.png}.

  The \textbf{Label Propagation} tab loads.

  \begin{figure}[H]
  \centering
  \includegraphics[width=0.6\textwidth]{images/vl-labelpropagation-tab.png}
  \caption{Label Propagation Tab}
  \end{figure}
\item
  To start, you need to select all of the class labels that are relevant
  for the dataset.

  Click
  
\includegraphics[width=0.04\linewidth,height=\textheight,keepaspectratio]{images/vl-labelpropagation-chooseclasses.png}

  The \textbf{Add Labels} dialog pops open:

  \begin{figure}[H]
  \centering
  \includegraphics[width=0.5\textwidth]{images/vl-labelpropagation-addlabels.png}
  \caption{Add Labels dialog}
  \end{figure}
\item
  Click in the field and choose labels from the dropdown list.

  This list contains the predefined classes from the
  \texttt{classes.json} file provided by your AOI machine and loaded
  during \hyperref[creating-datasets-from-machine-output]{dataset
  creation}. Select all classes that appear in your dataset because the
  system will attempt to label all data.

  \importantstart

  You must select \textbf{\emph{at least two classes}} for which you'd
  like to prepare labeling; classes not defined during this process
  cannot be assigned. \newline

  Consider your complete taxonomy before proceeding. \importantend
\item
  Click \textbf{Add}.

  The label classes appear in the right panel.
\item
  For each class label, select a minimum of 5 seed examples:

  \begin{itemize}
  \item
    \textbf{Drag and Drop} - For a single image or for a single cluster
    (in its entirety), drag it directly into the designated class area
    with visual feedback provided OR,
  \item
    \textbf{Select \& add} - Select all relevant images and then from
    the floating bar choose the relevant class and click \textbf{Assign
    Label}:
  \end{itemize}

  \begin{figure}
  \centering
  
\includegraphics[width=0.6\linewidth,height=\textheight,keepaspectratio,alt={Assign Label button}]{images/vl-labelprop-assignlabel.png}
  \caption{Assign Label button}
  \end{figure}

  Each class shows a preview of the examples that you assigned.

  \begin{figure}
  \centering
  \includegraphics[width=0.4\linewidth,height=\textheight,keepaspectratio,alt={Seed Examples preview}]{images/vl-labelpropagation-seed-examples-preview.png}
  \caption{Seed Examples preview}
  \end{figure}
\item
  Review and refine your selections:

  \begin{itemize}
  \item
    To remove an example, hover over it and click
    
\includegraphics[width=0.05\linewidth,height=\textheight,keepaspectratio,alt={Remove Example icon}]{images/vl-labelpropagation-deleteremove-labelclass-icon.png}.
  \item
    Progress indicators show completion toward minimum requirements:

    \begin{figure}[H]
    \centering
    \includegraphics[width=0.35\textwidth]{images/vl-labelpropagation-progress-indicator.png}
    \caption{Progress Indicator}
    \end{figure}
  \item
    To manage classes as needed during the seed creation process, click
    \textbf{Add Labels} to add additional classes for propagation.
  \item
    Click
    
\includegraphics[width=0.05\linewidth,height=\textheight,keepaspectratio,alt={Remove example icon}]{images/vl-labelpropagation-deleteremove-labelclass-icon.png}
    from the row with the class name to delete the entire sample set
    (and the label) from propagation.
  \end{itemize}
\item
  When ready, click \textbf{Submit}.

  \begin{figure}[H]
  \centering
  \includegraphics[width=0.4\textwidth]{images/vl-labelpropagation-runningspinner.png}
  \caption{Label Propagation Running}
  \end{figure}

  \textbf{Visual Layer} analyzes your seed examples and begins labeling
  images, creating sample sets for each of the class labels you
  configured.

  The system assigns confidence scores to each image for every possible
  label. Images with high confidence scores, typically 90\% or higher,
  are automatically labeled. Images with low confidence scores are sent
  to review batches for your input. When complete, a notification for
  review appears in the interface:

  \begin{figure}[H]
  \centering
  \includegraphics[width=0.35\textwidth]{images/vl-review-batch-notification.png}
  \caption{Review Ready Notification}
  \end{figure}
\end{enumerate}

\processstart

\textbf{Next Steps} \processend

When you receive notification, it's time to
\hyperref[reviewing-and-providing-feedback]{review} the labels added by
\textbf{Visual Layer}, and continue providing feedback to prepare the
system for automated propagation.

\begin{figure}
\centering
\includegraphics[width=0.4\linewidth,height=\textheight,keepaspectratio,alt={Progress and Review}]{images/vl-labelpropagation-progressbars.png}
\caption{Progress and Review}
\end{figure}

\chapter{Label Propagation: Iterative Review to
Completion}\label{reviewing-and-providing-feedback}

After submitting seeds, label propagation enters the iterative review
cycle. \textbf{Visual Layer} uses a confidence-based approach that
creates review batches when it's uncertain about image classifications.
This process ensures labeling accuracy while maintaining efficiency and
continues until the model reaches sufficient confidence across your
dataset.

\section{How Review Batches Work}\label{how-review-batches-work}

\textbf{Visual Layer} uses a confidence-based review process to ensure
labeling accuracy while maintaining efficiency. The algorithm assigns
confidence scores to each image for every possible class, creating
targeted review batches for human verification.

\renewcommand{\arraystretch}{1.5}
\rowcolors{1}{tableprimary}{tablesecondary}

\begin{tabularx}{\textwidth}{|>{\centering\arraybackslash}m{4em}|p{0.35\textwidth}|X|}
\hline
\textbf{Step} & \textbf{Phase} & \textbf{Description} \\
\hline

\raisebox{-7pt}{
\includegraphics[width=5.5mm]{images/1light.png}}
  & \textbf{Confidence Scoring}
  & After analyzing your seed examples, \textbf{Visual Layer} assigns confidence scores (0--100\%) to each image for every possible class. \\
\hline

\raisebox{-7pt}{
\includegraphics[width=5.5mm]{images/2light.png}}
  & \textbf{Automatic Labeling}
  & Images with high confidence scores (typically 90\% or higher) are automatically labeled and do not require review. \\
\hline

\raisebox{-7pt}{
\includegraphics[width=5.5mm]{images/3light.png}}
  & \textbf{Review Batch Creation}
  & Images with low confidence scores are sent to review batches. Each review round includes a maximum of 5--10\% of your total dataset. Example: with 1{,}000 images, you review approximately 50--100 images. \\
\hline

\raisebox{-7pt}{
\includegraphics[width=5.5mm]{images/4light.png}}
  & \textbf{Human Feedback}
  & Label propagation pauses and waits for your feedback. The process does not continue until the entire review batch is completed. Review automatically generated labels, especially items flagged for attention. You approve, correct, or reassign labels for uncertain images. Your feedback directly improves model accuracy in subsequent iterations. \\
\hline

\raisebox{-7pt}{
\includegraphics[width=5.5mm]{images/5light.png}}
  & \textbf{Model Refinement}
  & \textbf{Visual Layer} incorporates feedback and relabels the dataset with improved accuracy. This cycle repeats, generating new review batches as the model refines its understanding. \\
\hline

\end{tabularx}

\subsection{Source Indicators}\label{source-indicators}

\renewcommand{\arraystretch}{1.5}
\rowcolors{1}{tableprimary}{tablesecondary}
\begin{tabularx}{\textwidth}{|>{\centering\arraybackslash}m{1.2cm}|p{0.15\textwidth}|p{0.18\textwidth}|p{0.22\textwidth}|X|}
\hline
 & \textbf{Action} & \textbf{Use Case} & \textbf{What Happens} & \textbf{Important Notes} \\
\hline

\raisebox{-0.7\height}{
\includegraphics[width=.6cm]{images/vl-remove-redx-icon.png}}
  & \textbf{Incorrect (Reject)}
  & Certain it\textquotesingle s wrong class, unsure of which label is correct
  & Item moves to next-best predicted class. Remains flagged ``Requires Review''. Can continue rejecting until right class found
  & Undo available via toast notification. Example: Unknown defect types cycle through suggestions until recognized \\
\hline

\raisebox{-0.7\height}{
\includegraphics[width=.6cm]{images/vl-ignore-icon.png}}
  & \textbf{Ignore}
  & Outlier that doesn\textquotesingle t match any existing class
  & Item excluded from label propagation. Placed in ``Unlabeled'' cluster. Appears grayed out. Doesn\textquotesingle t count toward progress
  & System maintains 5--10\% review quota. Too many ignored items may trigger additional review batches \\
\hline

\raisebox{-0.7\height}{
\includegraphics[width=.6cm]{images/vl-approve-icon.png}}
  & \textbf{Approve}
  & Item correctly matches assigned class
  & Green border indicates approval. Advances to next item
  & Confirms model prediction \\
\hline

\raisebox{-0.7\height}{
\includegraphics[width=.6cm]{images/vl-labelpropagation-assigndifferentlabel-icon.png}}
  & \textbf{Assign to Other Class}
  & Assign to a different existing class
  & Manually reassigned to selected class. Advances to next item
  & Directly teaches model correct classification \\
\hline

\end{tabularx}

\tipstart

All VL and User items remain editable regardless of source. Previous
review decisions can be corrected if needed. \tipend

\importantstart

Items are organized by their suggested class based on highest confidence
score. \newline

You can provide feedback on both review items and high-confidence items.
\newline

All feedback helps improve subsequent labeling iterations. \importantend

\subsection{Feedback Actions}\label{feedback-actions}

During review rounds, you can provide feedback to the system via the
feedback actions:

\begin{figure}
\centering

\includegraphics[width=0.6\linewidth,height=\textheight,keepaspectratio,alt={Feedback Options}]{images/vl-labelpropagation-review-bar.png}
\caption{Feedback Options}
\end{figure}

\newpage

Following are the feedback actions and how to use each:

\renewcommand{\arraystretch}{1.5}
\rowcolors{1}{tableprimary}{tablesecondary}
\begin{tabularx}{\textwidth}{|>{\centering\arraybackslash}m{1.2cm}|p{0.15\textwidth}|p{0.18\textwidth}|p{0.22\textwidth}|X|}
\hline
 & \textbf{Action} & \textbf{Use Case} & \textbf{What Happens} & \textbf{Important Notes} \\
\hline

\raisebox{-0.7\height}{
\includegraphics[width=.6cm]{images/vl-remove-redx-icon.png}}
  & \textbf{Incorrect (Reject)}
  & Certain it's wrong class, unsure of which label is correct
  & Item moves to next-best predicted class\newline
    Remains flagged ``Requires Review''\newline
    Can continue rejecting until right class found
  & Undo available via toast notification\newline
    Example: Unknown defect types cycle through suggestions until recognized \\
\hline

\raisebox{-0.7\height}{
\includegraphics[width=.6cm]{images/vl-ignore-icon.png}}
  & \textbf{Ignore}
  & Outlier that doesn't match any existing class
  & Item excluded from label propagation\newline
    Placed in ``Unlabeled'' cluster\newline
    Appears grayed out\newline
    Does not count toward progress
  & System maintains 5–10\% review quota\newline
    Too many ignored items may trigger additional review batches \\
\hline

\raisebox{-0.7\height}{
\includegraphics[width=.6cm]{images/vl-approve-icon.png}}
  & \textbf{Approve}
  & Item correctly matches assigned class
  & Green border indicates approval\newline
    Advances to next item
  & Confirms model prediction \\
\hline

\raisebox{-0.7\height}{
\includegraphics[width=.6cm]{images/vl-labelpropagation-assigndifferentlabel-icon.png}}
  & \textbf{Assign to Other Class}
  & Assign to a different existing class
  & Manually reassigned to selected class\newline
    Advances to next item
  & Directly teaches model correct classification \\
\hline
\end{tabularx}

\processstart

\textbf{To start a review round} \processend

\begin{enumerate}
\def\labelenumi{\arabic{enumi}.}
\item
  When a review batch is ready for your review, a notification appears
  in the interface:

  \begin{figure}[H]
  \centering
  
\includegraphics[width=.5\textwidth]{images/vl-labelpropagation-needsreview-notification.png}
  \caption{Review Batch notification}
  \end{figure}

  The labeled images are automatically grouped by predicted class/label.
\item
  Click \textbf{Review} to begin reviewing images that were
  automatically labeled by \textbf{Visual Layer}.

  \notestart

  During review, all exploration features like filters, search, and
  visual similarity are disabled. \noteend

  \textbf{Visual Layer} automatically loads one of the clustered classes
  for review. The view displays all images that were assigned the same
  predicted label, grouped together for efficient review:

  \begin{figure}[H]
  \centering
  \includegraphics[width=0.9\textwidth]{images/vl-camtek-labelpropagation-reviewround-firstset.png}
  \caption{Review Round - First Class Loaded}
  \end{figure}

  You can search for specific classes if you have many classes to
  review:

  \begin{figure}[H]
  \centering
  \includegraphics[width=0.5\textwidth]{images/vl-review-search-classes.png}
  \caption{Search classes}
  \end{figure}

  Within each class, you can switch between two tabs:

  \begin{itemize}
  \item
    \textbf{For Review} (default tab) - Displays images that need your
    review, marked with
    
\includegraphics[width=0.15\linewidth,height=\textheight,keepaspectratio,alt={Requires Review badge}]{images/vl-labelprop-label-forreview.png}.
    Once you review an item, it's marked with
    
\includegraphics[width=0.15\linewidth,height=\textheight,keepaspectratio,alt={Reviewed badge}]{images/vl-labelprop-label-reviewed.png}.
  \item
    \textbf{Annotated} - Displays seed samples you provided and images
    that \textbf{Visual Layer} labeled with high confidence (typically
    ≥90\%). All items are marked with
    
\includegraphics[width=0.13\linewidth,height=\textheight,keepaspectratio,alt={Annotated badge}]{images/annotated-badge.png}
    and show their source indicator (VL for automated labels, User for
    reviewed items, or Seed for original examples).
  \end{itemize}

  \begin{figure}[H]
  \centering
  \includegraphics[width=0.9\textwidth]{images/vl-camtek-labelpropagation-reviewround-firstset-annotated.png}
  \caption{Annotated Tab View}
  \end{figure}

  You can also view all clusters included in the propagation, similar to
  other cluster views in \textbf{Visual Layer}. Each cluster card
  displays:

  \begin{itemize}
  \tightlist
  \item
    Total number of images in the cluster for this review round
    (including both Annotated and For Review images)
  \item
    Number of For Review images specifically for that cluster
  \end{itemize}

  \begin{figure}[H]
  \centering
  \includegraphics[width=0.9\textwidth]{images/vl-camtek-labelpropagation-reviewround-allclusters.png}
  \caption{All Clusters View}
  \end{figure}

  \notestart

  If a class has no images for review in the current round, its cluster
  card appears but shows no content, as seen with ``Bump Scratch'' in
  the example above. \noteend
\item
  Review every image for every class:

  \begin{itemize}
  \item
    Provide feedback directly from the default display.
  \item
    Alternatively, double-click any image to open a modal carousel view.

    \begin{figure}[H]
    \centering
    \includegraphics[width=0.7\textwidth]{images/vl-review-carousel-modal.png}
    \caption{Carousel Modal}
    \end{figure}
  \item
    Navigate through individual items using carousel controls.
  \item
    Zoom in/out for detailed defect inspection.
  \end{itemize}
\item
  For every automatically labeled image, provide the model with
  feedback:

  \begin{itemize}
  \tightlist
  \item
    Approve
  \item
    Reject
  \item
    Ignore
  \item
    Apply different label
  \end{itemize}

  See \hyperref[feedback-actions]{Feedback Actions} for more
  information.
\item
  Review each item in every class until all items are processed.
\item
  Click \textbf{Submit} to send all your feedback to the system. Label
  propagation continues with the next iteration, using your feedback to
  improve accuracy in subsequent labeling cycles. When the system has
  high confidence, it automatically applies labels during the relevant
  cycle.
\item
  Once the system is confident and applies all labels across your entire
  dataset, a notification pops up, similar to the following:

  \begin{figure}[H]
  \centering
  \includegraphics[width=0.6\textwidth]{images/vl-labelpropagation-complete-notification.png}
  \caption{Propagation Success}
  \end{figure}
\end{enumerate}

\section{Monitoring Progress}\label{monitoring-progress}

Throughout the label propagation process, \textbf{Visual Layer} provides
multiple progress indicators to help you understand completion status
and make informed decisions about when to finish.

\processstart

\textbf{Training readiness tracking} \processend

\textbf{Visual Layer} requires 100 labeled images per class as the
minimum threshold for training readiness. The platform provides three
distinct progress indicators to track your progress toward this goal and
overall labeling completion:

\begin{itemize}
\item
  \textbf{Train-worthy Progress} -- Shows progress toward the 100 images
  per class training readiness goal. When this reaches 100\%, you have
  met the minimum threshold required to train models.

  \begin{figure}[H]
   \centering
   \includegraphics[width=0.5\textwidth]{images/vl-review-trainworthy-progress.png}
   \caption{Train-worthy Progress}
   \end{figure}
\item
  \textbf{Overall Progress} -- Shows the percentage of your entire
  dataset that has been labeled, regardless of the per-class training
  threshold. This indicates how much of your total dataset has received
  labels.

  \begin{figure}[H]
   \centering
   \includegraphics[width=0.5\textwidth]{images/vl-review-overall-progress.png}
   \caption{Overall Progress}
   \end{figure}
\item
  \textbf{Per-Class Progress} -- Shows how many labeled images each
  individual class has toward the 100-image minimum threshold. Use this
  to identify which classes need more labeled examples.

  \begin{figure}[H]
   \centering
   \includegraphics[width=0.5\textwidth]{images/vl-review-perclass-progress.png}
   \caption{Per-Class Progress}
   \end{figure}
\end{itemize}

\section{Managing Iterations and
Completion}\label{managing-iterations-and-completion}

This section outlines how to manage the lifecycle of a label propagation
run after seed creation and initial labeling:

\begin{itemize}
\tightlist
\item
  finishing the run early to finalize labels
\item
  resetting the process if the results are unsatisfactory
\end{itemize}

\subsection{Finish Label Propagation
Early}\label{finish-label-propagation-early}

Once the system is confident and applies all labels across your entire
dataset, a notification pops up.

You can finish the process early, enabling Visual Layer to apply labels
across your dataset before completing all iterations. This is available
when \textbf{Needs Review} status is active only.

\begin{figure}
\centering
\includegraphics[width=0.35\linewidth,height=\textheight,keepaspectratio,alt={Finish Early}]{images/vl-labelpropagation-review-finish-panel.png}
\caption{Finish Early}
\end{figure}

From the review details in the right panel:

\begin{enumerate}
\def\labelenumi{\arabic{enumi}.}
\item
  Click \textbf{Finish} when satisfied with current progress.
\item
  Confirm the action in the dialog that appears.

  The system:

  \begin{itemize}
  \tightlist
  \item
    Stops label propagation immediately
  \item
    Writes all labeled data permanently to your dataset
  \item
    Leaves remaining unlabeled data unchanged
  \end{itemize}
\end{enumerate}

\subsection{Resetting the Process}\label{resetting-the-process}

You can reset the process only at stages when \textbf{Review} is
available. Reset \textbf{\textbf{permanently deletes}} all label
propagation progress and cannot be undone. Your original dataset and any
labels from previous completed label propagation runs remain unchanged.

Reset label propagation if you:

\begin{itemize}
\tightlist
\item
  notice consistent labeling errors that suggest poor seed examples
\item
  want to try a different class taxonomy or seed strategy
\item
  the results are not meeting your accuracy expectations
\item
  accidentally started with incorrect seed examples
\end{itemize}

\processstart

\textbf{To reset the process} \processend

\criticalstart

Reset \textbf{\textbf{permanently deletes}} all label propagation
progress and cannot be undone. Your original dataset and any labels from
previous completed label propagation runs remain unchanged. \criticalend

\begin{enumerate}
\def\labelenumi{\arabic{enumi}.}
\tightlist
\item
  From the right panel, click
  
\includegraphics[width=0.08\linewidth,height=\textheight,keepaspectratio]{images/vl-reset-button.png}.
\item
  Confirm the reset action in the warning dialog.
\item
  The system completely deletes all:

  \begin{itemize}
  \tightlist
  \item
    automatically labeled images (high-confidence assignments)
  \item
    user review feedback and corrections
  \item
    progress and iteration history
  \end{itemize}
\item
  Now, from the seed creation step you can:

  \begin{itemize}
  \tightlist
  \item
    Start fresh
  \item
    Modify your seed examples
  \item
    Add or remove classes if supported
  \end{itemize}
\end{enumerate}

\section{Next Steps}\label{next-steps}

When your dataset reaches train worthy status:

\begin{itemize}
\tightlist
\item
  \hyperref[evaluating-label-propagation]{Evaluate labeling results}
\item
  \hyperref[train]{Train your downstream models}
\end{itemize}

\genericdivider{Step 3: Train Models Downstream}

\chapter{Train Models: How It Works}\label{train-models-how-it-works}

Training models downstream is the third and final main step in the
\textbf{Visual Layer} workflow. Once label propagation reaches
\hyperref[monitoring-progress]{train-worthy status} and you have
\hyperref[evaluating-label-propagation]{validated results}, you can send
your finalized, validated dataset to your model-training pipeline for
deployment.

You can train a model downstream when the dataset has achieved
\hyperref[monitoring-progress]{\textbf{Train-Worthy} status,} allowing
you to send the complete dataset downstream to train your Camtek models.
This integrated training workflow eliminates manual data export and
preparation steps.

\textbf{Visual Layer} automatically formats your labeled dataset and
sends it to the Camtek training infrastructure with a single click. When
sending the data for training, the system uses all labeled data from
your completed label propagation, including seed examples, reviewed
items, and high-confidence automatic labels. No manual data preparation
is required. You can use a single dataset to support the training of
multiple models over time; however, you can only run one training
session per dataset at a time.

The downstream model training process consists of the following
substeps:

\begin{figure}[H]
\centering
\includegraphics[width=0.9\textwidth]{images/vl-train-models-diagram.png}
\caption{Train Models Downstream Workflow}
\end{figure}

\begin{enumerate}
\def\labelenumi{\arabic{enumi}.}
\tightlist
\item
  \textbf{Initiate Training} - Select your target Camtek model
  configuration and click Train to package and send your labeled dataset
\item
  \textbf{Monitor Progress} - Track training job status through the
  \textbf{Visual Layer} interface
\item
  \textbf{Review Results} - Access completed models in the Model Catalog
  to view training results and performance metrics
\end{enumerate}

\chapter{Train Models: Initiate}\label{train}

\processstart

\textbf{To initiate model training} \processend

\begin{enumerate}
\def\labelenumi{\arabic{enumi}.}
\item
  Navigate to the relevant dataset.
\item
  Click
  
\includegraphics[width=0.15\linewidth,height=\textheight,keepaspectratio]{images/vl-train-model-button.png}
  in the top right corner. The \textbf{New Auto Training Job} dialog
  opens, similar to the following:

  \begin{figure}[H]
  \centering
  \includegraphics[width=0.4\textwidth]{images/vl-labelpropagation-new-training-downstream.png}
  \caption{Training Job dialog}
  \end{figure}

  Following are details of the fields in this window:
\end{enumerate}

\begin{longtable}[]{@{}
  >{\raggedright\arraybackslash}p{(\linewidth - 4\tabcolsep) * \real{0.3333}}
  >{\raggedright\arraybackslash}p{(\linewidth - 4\tabcolsep) * \real{0.3333}}
  >{\raggedright\arraybackslash}p{(\linewidth - 4\tabcolsep) * \real{0.3333}}@{}}
\toprule\noalign{}
\begin{minipage}[b]{\linewidth}\raggedright
\textbf{Component}
\end{minipage} & \begin{minipage}[b]{\linewidth}\raggedright
\textbf{Description}
\end{minipage} & \begin{minipage}[b]{\linewidth}\raggedright
\textbf{Example/Notes}
\end{minipage} \\
\midrule\noalign{}
\endhead
\bottomrule\noalign{}
\endlastfoot
\textbf{Dataset Name} & The human-readable name of the \textbf{Visual
Layer} dataset being sent for training. & Populated automatically from
the current dataset. (Read-only, for reference); you can change the name
from the title bar inside the dataset itself. \\
\textbf{Dataset ID} & The unique system identifier for the dataset. &
For example: 0f4698-9940-11f0-8218-6ad3b2ffe3e5 (Read-only, for
reference) \\
\textbf{Model Name} & Allows selection of the target Camtek model
configuration for this training run. & Select model from the existing
list of configured Camtek models. \\
\textbf{Train button} & Initiates the training job. & Click to package
the labeled data and send it to the Camtek training infrastructure. \\
\end{longtable}

\begin{enumerate}
\def\labelenumi{\arabic{enumi}.}
\setcounter{enumi}{2}
\item
  From \textbf{Model Name} choose the relevant model.
\item
  Click \textbf{Train}.

  \begin{figure}[H]
  \centering
  \includegraphics[width=0.4\textwidth]{images/vl-labelprop-trainmodel-datasetsentsuccessfully.png}
  \caption{Training Success}
  \end{figure}

  \textbf{Visual Layer} automatically packages your labeled dataset with
  all metadata and sends it to your training infrastructure.

  Optionally, you can view the training progress, which now appears in
  the action bar:

  \begin{figure}[H]
  \centering
  
\includegraphics[width=0.4\textwidth]{images/vl-trainmodel-progress.png}
  \caption{Training Progress indicator}
  \end{figure}

  Training typically completes within several hours depending on dataset
  size and model complexity. Once a model training job ends
  successfully, you'll receive a notification with a direct action to
  view the result in the \hyperref[the-model-catalog]{\textbf{Model
  Catalog}}.
\end{enumerate}

\chapter{Train Models: Model Catalog}\label{the-model-catalog}

The \textbf{Model Catalog} is the central hub for managing your
downstream model training and results. Once an external model training
job successfully completes, the trained model data is accessible from
the \textbf{Model Catalog}. From here, you can:

\begin{itemize}
\tightlist
\item
  View the detailed training results
\item
  Copy a link to the results report path
\item
  Delete the model from the \textbf{Visual Layer} catalog
\end{itemize}

\processstart

\textbf{To get started} \processend

From the left of the \textbf{Visual Layer} interface, click
\includegraphics[width=0.06\linewidth,height=\textheight,keepaspectratio,alt={Catalog icon}]{images/vl-modelcatalog-icon.png}.

The \textbf{Model Catalog} loads, and appears similar to the following:

\section{Catalog Structure and
Features}\label{catalog-structure-and-features}

The \textbf{Model Catalog} displays a table view of all the custom
models you've sent for training. You can view and filter the table based
on your needs.

The table displays the following information:

\begin{longtable}[]{@{}
  >{\raggedright\arraybackslash}p{(\linewidth - 2\tabcolsep) * \real{0.5000}}
  >{\raggedright\arraybackslash}p{(\linewidth - 2\tabcolsep) * \real{0.5000}}@{}}
\toprule\noalign{}
\begin{minipage}[b]{\linewidth}\raggedright
\textbf{Column}
\end{minipage} & \begin{minipage}[b]{\linewidth}\raggedright
\textbf{Description}
\end{minipage} \\
\midrule\noalign{}
\endhead
\bottomrule\noalign{}
\endlastfoot
\textbf{Model Name} & The name of the selected model, as returned from
Camtek. \\
\textbf{Type} & The type of the model. This is currently always
\textbf{Classification}. \\
\textbf{Source} & The training source. This is currently always Camtek
Auto ML. \\
\textbf{Dataset Name} & The human-readable name of the dataset in
\textbf{Visual Layer} used for training. \\
\textbf{Dataset ID} & The unique system identifier of the training
dataset in \textbf{Visual Layer}. \\
\textbf{Job} & The production job identifier for the training job. \\
\textbf{Setup} & The machine setup configuration used. \\
\textbf{Recipe List} & A list of the scanning recipes applied. \\
\textbf{Tool List} & The tools used. \\
\textbf{Last Updated} & The timestamp from the training job results
folder. \\
\textbf{Results} & An icon that opens the HTML report for the model
results. \\
\textbf{Production} & Check or uncheck based on your training
processes. \\
\end{longtable}

\section{Model Actions}\label{model-actions}

The three-dot menu provides key functions for managing the trained
models.

\subsection{View Results (HTML)}\label{view-results-html}

From the relevant row in the \textbf{Model Catalog},click the the
three-dot menu and choose \textbf{Export}.

The entire model results folder downloads as a ZIP file to the user's
local `Downloads' folder.

\subsection{Delete a Model}\label{delete-a-model}

Users with the necessary permissions can remove any model from the
display in the \textbf{Visual Layer} catalog.

\processstart

\textbf{To delete a model from the catalog} \processend

\begin{enumerate}
\def\labelenumi{\arabic{enumi}.}
\item
  From the relevant row in the \textbf{Model Catalog},click the the
  three-dot menu and choose \textbf{Delete}.

  A confirmation modal appears.
\item
  Confirm deletion (or cancel to stop).

  The model entry is immediately removed from the catalog interface.
\end{enumerate}

\sectiondivider{Advanced Dataset Features}

\chapter{Adding Media to Datasets}\label{adding-media-to-datasets}

As production continues, inspection machines continuously generate new
data. New images may contain defects you're already tracking, introduce
entirely new defect types, or present a combination of both.

You can add these images to existing datasets without creating new ones.
This maintains dataset continuity across production runs and preserves
your existing class taxonomy. After re-indexing, you can re-run label
propagation on the complete dataset, reusing your proven seed examples
while Visual Layer learns from both existing and newly added images
together.

Adding media to an existing dataset enables you to:

\begin{itemize}
\tightlist
\item
  Expand training datasets with new production examples without
  fragmenting data across multiple datasets
\item
  Maintain consistent Job Recipe and class taxonomy across all images in
  a dataset
\item
  Continue \hyperref[re-indexing-the-updated-dataset]{label propagation}
  workflows on newly added images using existing seed examples and model
  learning
\item
  Track dataset evolution over time while preserving the original
  dataset ID and configuration
\end{itemize}

\section{How It Works}\label{how-it-works}

\textbf{Visual Layer} uses a two-stage process when adding media to
optimize performance and give you control over when computationally
intensive operations run. New media becomes immediately visible for
\hyperref[adding-media-to-an-existing-dataset]{exploration} after
upload. You'll trigger indexing \textbf{separately} when you're ready.

\newpage

\textbf{Here's how it works:}

\renewcommand{\arraystretch}{1.5}
\rowcolors{1}{tableprimary}{tablesecondary}

\begin{tabularx}{\textwidth}{|>{\centering\arraybackslash}m{4em}|p{0.35\textwidth}|X|}
\hline
\textbf{Step} & \textbf{Phase} & \textbf{Description} \\
\hline
\raisebox{-7pt}{
\includegraphics[width=5.5mm]{images/1light.png}}
  & \textbf{Upload and Validate}
  & Upload new AOI output folders. \textbf{Visual Layer} validates that the metadata matches the existing dataset exactly. \\
\hline
\raisebox{-7pt}{
\includegraphics[width=5.5mm]{images/2light.png}}
  & \textbf{Partial Update}
  & \textbf{Visual Layer} automatically generates embeddings for visual search, creates thumbnails, processes metadata, and groups new images into a temporary "New Media" cluster visible immediately for exploration. This cluster is part of your existing dataset, but separate from all other clusters during this stage, enabling you to review the added media before re-indexing completely. Additionally, you can repeat these steps to continue adding additional folders if you didn't include all relevant media initially. When this update is complete, all media that was newly added is in \textbf{Pending Indexed} status.\\
\hline
~ &  &  \\
\hline
\raisebox{-7pt}{
\includegraphics[width=5.5mm]{images/3light.png}}
  & \textbf{Manually Trigger Re-indexing}
  & When ready, you can trigger full indexing. During indexing, \textbf{Label Propagation} is unavailable. \textbf{Visual Layer} re-clusters all images (new and existing), regenerates quality issue detection, and updates cluster IDs. Dataset status changes to \textbf{Indexing}. \\
\hline
\raisebox{-7pt}{
\includegraphics[width=5.5mm]{images/4light.png}}
  & \textbf{Ready for Use}
  & Indexing completes and dataset status returns to \textbf{Ready} status. All media in the datset is set to \textbf{Indexed} status. New images are fully integrated in the dataset, with all clusters refreshed accordingly, and available for label propagation and training. \\
\hline
\end{tabularx}

\notestart

Dataset ID and all image IDs remain unchanged during indexing. Only
cluster IDs are updated to reflect the new clustering results. \noteend

\section{Adding Media to an Existing
Dataset}\label{adding-media-to-an-existing-dataset}

Visual Layer uses a two-stage process when adding media to optimize
performance and give you control over when computationally intensive
operations run. New media becomes immediately visible for
\hyperref[adding-media-to-an-existing-dataset]{exploration} after
upload, with full indexing triggered manually when you're ready. You can
upload multiple batches and explore before triggering the clustering and
analysis steps.

\criticalstart

New media must match the metadata of the original dataset. Mismatched
metadata will cause validation failures and prevent media from being
added. \newline

To work with datasets that have different metadata,
\hyperref[creating-datasets-from-machine-output]{create a new dataset}
instead. \criticalend

\notestart

See \hyperref[validation-requirements]{Dataset and Media Validation
Requirements} in the Appendices for detailed validation requirements,
file structure specifications, and common errors when adding media to
existing datasets. \noteend

\processstart

\textbf{To add media to an existing dataset} \processend

\begin{enumerate}
\def\labelenumi{\arabic{enumi}.}
\item
  Navigate to the dataset where you want to add new images.
\item
  Go to the \textbf{Data} tab.

  The \textbf{Data} area loads, similar to the following:

  \begin{figure}[H]
   \centering
   \includegraphics[width=0.8\textwidth]{images/vl-dataset-updatedataset.png}
   \caption{Data Tab for Adding Media}
   \end{figure}

  This tab uses the same interface as
  \hyperref[creating-datasets-from-machine-output]{dataset creation},
  but adds media to the current dataset instead of creating a new one.
  The original folder/s you used to create the dataset appear in the
  \textbf{Select Folders} table, and you can add additional folders as
  well.

  The name of the dataset you're updating appears at the top left
  accordingly. Hover over the
  
\includegraphics[width=5mm]{images/vl-info-i-icon.png} next to the
  dataset name to confirm its details:

  \begin{figure}[H]
   \centering
   \includegraphics[width=0.3\textwidth]{images/vl-dataset-updatedataset-i-icon-hover.png}
   \caption{Dataset Details Tooltip}
   \end{figure}
\item
  Click \textbf{Browse Folders} and select your AOI output folder
  containing the new images and metadata files.

  \notestart

  You can add multiple folders to a single dataset, as described in the
  next step. \noteend
\item
  \textbf{Visual Layer} extracts metadata and validates folder structure
  and contents.

  When complete, the folder appears in a table on the same screen.
\item
  Upload additional folders that share identical metadata with the
  original dataset.

  \tipstart If you have multiple production runs to add, upload them all
  before triggering indexing. This reduces processing time and ensures
  the clustering algorithm analyzes your complete dataset together.
  \tipend

  \textbf{Visual Layer} validates Job/Setup/Recipe compatibility against
  the existing dataset to ensure all folders originate from the same
  production run.

  \begin{figure}
  \centering
  \includegraphics[width=0.8\linewidth,height=\textheight,keepaspectratio,alt={Selected Folders table}]{images/vl-create-dataset-selectfolders.png}
  \caption{Selected Folders table}
  \end{figure}
\item
  Folders appear in the table showing their compatibility status:

  \notestart

  See \hyperref[creating-datasets-from-machine-output]{Creating a
  Dataset} for additional details. \noteend
\item
  Filter the \textbf{Selected Folders} table to identify compatible
  folders:

  Since you may not know the folder origins in advance, upload as many
  folders as relevant and then use the table filters to identify those
  that share identical Job/Setup/Recipe metadata with the existing
  dataset.

  \begin{itemize}
  \item
    \textbf{Filter by production metadata} - Use the Job Name, Setup
    Name, and Recipe Name filters to identify folders from the same
    production run:

    \includegraphics[width=0.7\linewidth,height=\textheight,keepaspectratio]{images/vl-create-dataset-selectfolders-filters.png}
  \item
    \textbf{Check compatibility} - Hover over the info icon to view
    details about any incompatible folders:

    \includegraphics[width=0.7\linewidth,height=\textheight,keepaspectratio]{images/vl-create-dataset-selectfolders-error.png}
  \end{itemize}

  Compatible folders remain selectable, while incompatible folders
  (those with different metadata) are grayed out and cannot be selected.
\item
  Mark the folders from the list that should be added to the dataset and
  click \textbf{Continue}.

  \begin{figure}[H]
  \centering
  \includegraphics[width=0.9\textwidth]{images/vl-dataset-updatedataset-selectfolders-alreadyexists.png}
  \caption{Selected Folders for Adding Media}
  \end{figure}

  \textbf{Visual Layer} loads a preview of the intended dataset update
  for you to double-check and approve.
\item
  Click \textbf{Add Media}.

  \begin{figure}[H]
  \centering
  \includegraphics[width=0.6\textwidth]{images/vl-create-dataset-progress.png}
  \caption{Upload Progress steps}
  \end{figure}
\item
  When upload completes, a notification confirms that new media was
  successfully added. All newly added media is marked with status
  \textbf{Pending Indexed}, for which you can filter when exploring your
  dataset:
\end{enumerate}

\begin{figure}
\centering
\pandocbounded{\includegraphics[keepaspectratio,alt={Pending Indexed}]{images/vl-addmedia-pendingindexed.png}}
\caption{Pending Indexed}
\end{figure}

\section{Re-indexing the Updated
Dataset}\label{re-indexing-the-updated-dataset}

After reviewing your new clusters, you can re-index the entire dataset
in order to commence a new
\hyperref[re-indexing-the-updated-dataset]{label propagation} round.

\begin{enumerate}
\def\labelenumi{\arabic{enumi}.}
\item
  When ready to proceed with
  \hyperref[re-indexing-the-updated-dataset]{label propagation} or
  \hyperref[train]{training}, click \textbf{Re-Index Dataset}. If you
  plan to add more media, wait and batch multiple uploads before
  indexing.

  Full indexing begins. The dataset status changes to \textbf{Indexing},
  and the following features become temporarily unavailable:
  \textbf{Label Propagation} and \textbf{Adding Media}.

  \notestart

  Track indexing progress in the
  \textbf{\hyperref[task-manager-statuses]{Task Manager}}, which shows
  task type, dataset name and ID, duration, and status. \noteend
\item
  When indexing finishes, a notification appears confirming the dataset
  is ready.

  All features are re-enabled. Your newly added images are now fully
  integrated and available for
  \hyperref[re-indexing-the-updated-dataset]{label propagation} and
  \hyperref[train]{training}.
\end{enumerate}

\section{Updating Dataset Labels}\label{updating-dataset-labels}

Once the dataset has been reindexed, you can return to
\hyperref[re-indexing-the-updated-dataset]{label propagation}. Assuming
the labels you initially used are still relevant for images in the
dataset, we recommend you approach this as follows:

\begin{enumerate}
\def\labelenumi{\arabic{enumi}.}
\item
  From inside the relevant dataset, initiate
  \hyperref[re-indexing-the-updated-dataset]{label propagation}, adding
  all of the relevant class labels.

  This includes any new labels that may now be relevant!
\item
  Filter for all of your original seed images for one of the relevant
  labels.
\item
  Select one, and then choose \textbf{Select All}.
\item
  Add them to the relevant label in the right panel.
\item
  Repeat the process for all other labels that you already identified
  seed examples for.
\end{enumerate}

\chapter{Evaluating Label
Propagation}\label{evaluating-label-propagation}

After completing label propagation, you can validate the accuracy of the
automatically generated labels. The \textbf{LP Auto Annotation
Validation Report} is an interactive HTML report comparing your Ground
Truth data with the labels generated through the label propagation
process.

Use the evaluation report to make informed decisions about your label
propagation quality:

\begin{itemize}
\tightlist
\item
  Measure overall accuracy to determine if your dataset is ready for
  model training
\item
  Identify which defect classes need additional seed examples or clearer
  definitions based on per-class precision
\item
  Track error sources by label origin to determine whether additional
  review iterations or annotation guideline clarification is needed
\item
  Locate classes with high missed detection rates requiring more diverse
  seed examples
\item
  Export validation data for compliance documentation and process audits
\end{itemize}

\section{Prerequisites}\label{prerequisites}

Before generating an evaluation report, ensure you have:

\begin{itemize}
\tightlist
\item
  Completed the \hyperref[how-label-propagation-works]{label
  propagation} feedback loop for the dataset
\item
  A Ground Truth file (\texttt{db.csv}) containing validated annotations
  for your images
\item
  A class mapping file (\texttt{classes.json}) defining the taxonomy
  used in your Ground Truth data
\end{itemize}

\notestart

The Ground Truth taxonomy must align with the classes used during label
propagation. Additionally, in this file, the Image Name column which
includes the relative path must include the folder name in which the
image is saved. \noteend

\section{Generate an Evaluation
Report}\label{generate-an-evaluation-report}

The report includes:

\begin{itemize}
\tightlist
\item
  High-level accuracy metrics and summary statistics
\item
  A breakdown of the user input accuracy
\item
  Confusion matrix showing classification performance across all classes
\item
  Image-by-image comparison of Ground Truth versus Label Propagation
  results
\item
  Visual examples highlighting agreement and discrepancies
\end{itemize}

\notestart

See \hyperref[report-structure-and-details]{Report Structure and
Details} for more information. \noteend

An example evaluation report appears similar to the following:

\begin{figure}[H]
\centering
\includegraphics[width=0.9\textwidth]{images/vl-camtek-evaluation-report-example.png}
\caption{Evaluation Report Example}
\end{figure}

\processstart

\textbf{To create an evaluation report} \processend

\begin{enumerate}
\def\labelenumi{\arabic{enumi}.}
\item
  Navigate to your dataset in \textbf{Visual Layer}.

  Ensure label propagation is complete for the dataset you want to
  evaluate.
\item
  Go to the three-dot menu and click
  \includegraphics[width=5mm]{images/vl-evaluation-icon.png}
  \textbf{Evaluate Label Propagation}.

  \notestart

  The \textbf{Evaluate Label Propagation} option appears only after
  label propagation has finished running. \noteend

  The \textbf{Evaluate Label Propagation} dialog appears.

  \includegraphics[width=0.5\linewidth,height=\textheight,keepaspectratio]{images/vl-labelpropagation-evaluation.png}
\item
  Click \textbf{Browse Folders} and provide the paths to your Ground
  Truth files:

  \begin{itemize}
  \tightlist
  \item
    \texttt{db.csv} - Your validated image classifications
  \item
    \texttt{classes.json} - The defect taxonomy definition
  \end{itemize}

  \importantstart

  The Ground Truth files must use the same class names and structure as
  your Label Propagation configuration. Mismatched taxonomies will
  prevent report generation. \importantend
\item
  Click \textbf{Run Evaluation}.

  A progress notification appears showing the comparison status. The
  system validates your files, matches image names, compares
  classifications, and generates the HTML report.
\item
  When processing completes, a notification confirms the report is
  ready.

  The report downloads automatically as a ZIP file to your Downloads
  folder. The ZIP contains:

  \begin{itemize}
  \tightlist
  \item
    \texttt{index.html} - located in the root of the folder , this is
    the report itself (viewable in Chrome)
  \item
    \texttt{js/} - Required JavaScript files (Chart.js and report
    functionality)
  \item
    \texttt{style/} - CSS styling files
  \item
    \texttt{DB.csv} - Merged comparison data (including class IDs)
  \item
    \texttt{classes.json} - Class mapping reference (maps IDs to class
    names)
  \item
    \texttt{missed\_images\_data.js} - Data file for mismatched samples
  \end{itemize}
\item
  Extract the ZIP file.
\item
  Double click the \texttt{index.html} file in the root of the folder to
  open the report in your web browser.
\end{enumerate}

\tipstart

Save the extracted report folder for your records. The report remains
fully functional offline. \tipend

\processstart

\textbf{To explore and analyze the report data} \processend

\begin{itemize}
\tightlist
\item
  In any table, hover over a cell or row to highlight its position more
  clearly.
\item
  Slice and dice the
  \hyperref[image-classification-validation]{\textbf{Image
  Classification Validation}} table via various filters.
\end{itemize}

\section{Report Structure and
Details}\label{report-structure-and-details}

The evaluation report contains the details and insights as described in
this section.

\subsection{General Information}\label{general-information}

The report header provides an at-a-glance understanding of how well your
label propagation performed across the entire dataset. This summary
displays key metrics and metadata about your validation:

\begin{figure}
\centering
\pandocbounded{\includegraphics[keepaspectratio,alt={General Informaion Example}]{images/vl-camtek-evaluationreport-sections1-2-examples.png}}
\caption{General Informaion Example}
\end{figure}

\begin{longtable}[]{@{}
  >{\raggedright\arraybackslash}p{(\linewidth - 2\tabcolsep) * \real{0.5000}}
  >{\raggedright\arraybackslash}p{(\linewidth - 2\tabcolsep) * \real{0.5000}}@{}}
\toprule\noalign{}
\begin{minipage}[b]{\linewidth}\raggedright
\textbf{Field}
\end{minipage} & \begin{minipage}[b]{\linewidth}\raggedright
\textbf{Description}
\end{minipage} \\
\midrule\noalign{}
\endhead
\bottomrule\noalign{}
\endlastfoot
\textbf{Generated On} & Date and time when the report was created
(YYYY-MM-DD HH:MM:SS format) \\
\textbf{Total Samples} & The total number of sample images that were
given for the seed from this dataset \\
\textbf{Number of Classes} & Total number of classes (labeled) included
in the label propagation \\
\textbf{Class breakdown} & A list of the classes, their class IDs, and
the total number of images in each class. For example, 1 - butterfly
(92) means the class ID is 1, the class name is butterfly, and 92 images
were labeled with this class. \\
\end{longtable}

\subsection{LP Auto Annotation
Performance}\label{lp-auto-annotation-performance}

This section evaluates how closely the Label Propagation results match
the Ground Truth, appearing similar to the following example:

\begin{figure}
\centering
\pandocbounded{\includegraphics[keepaspectratio,alt={LP Auto Annotation Performance Example}]{images/vl-camtek-evaluationreport-lpautoannotationperformance-examples.png}}
\caption{LP Auto Annotation Performance Example}
\end{figure}

\subsubsection{Per-Class Performance}\label{per-class-performance}

This table displays detailed classification performance metrics for each
class.

\rowcolors{1}{tableprimary}{tablesecondary}
\begin{tabularx}{\textwidth}{|p{0.25\textwidth}|X|p{0.3\textwidth}|}
\hline
\textbf{Metric} & \textbf{Description} & \textbf{Formula} \\
\hline
\textbf{Class} & The name of the class for these statistics &  \\
\hline
\textbf{Precision} & Percentage of Label Propagation predictions that were correct for this class & True Positives / (True Positives + False Positives) \\
\hline
\textbf{Recall} & Percentage of actual instances that Label Propagation correctly identified & True Positives / (True Positives + False Negatives) \\
\hline
\textbf{F1-Score} & Balanced measure combining Precision and Recall & 2 × (Precision × Recall) / (Precision + Recall) \\
\hline
\end{tabularx}

\subsubsection{Confusion Matrix}\label{confusion-matrix}

The confusion matrix table displays the predicted classes versus the
ground truth. Rows represent the actual Ground Truth labels, while
columns show the predicted Label Propagation results.

Each cell in the matrix displays the count of images where a Ground
Truth class (row) was predicted as a specific Label Propagation class
(column). The diagonal cells (where row equals column) represent correct
classifications.

The matrix includes:

\begin{itemize}
\tightlist
\item
  \textbf{Row Totals (GT Total)} - The total number of images for each
  Ground Truth class
\item
  \textbf{Column Totals (Predicted Total)} - The total number of images
  predicted for each class by Label Propagation
\end{itemize}

\begin{figure}
\centering
\pandocbounded{\includegraphics[keepaspectratio,alt={Confusion Matrix Example}]{images/vl-camtek-evaluation-confusionmatrix-example.png}}
\caption{Confusion Matrix Example}
\end{figure}

\notestart

For datasets with many classes, the matrix displays class IDs with a
mapping reference at the top of the table. Use horizontal scrolling to
view all columns. \noteend

\subsection{User Performance Overview}\label{user-performance-overview}

This section of the report provides an overview of the user's
performance during Label Propagation as compared to the ground truth,
which is divided into sub-sections:

\begin{figure}
\centering
\pandocbounded{\includegraphics[keepaspectratio,alt={User Performance Overview Example}]{images/vl-camtek-evaluation-userperformance-example.png}}
\caption{User Performance Overview Example}
\end{figure}

\subsubsection{Summary}\label{summary}

The summary itemizes:

\begin{itemize}
\tightlist
\item
  Total Questions Answered - the total number of feedback requests that
  the user responded to
\item
  Total User Errors - the number of errors that the user made in total
\item
  Seed-Stage Errors - the number of initial seed samples that the user
  provided incorrectly
\item
  Review-Stage Errors - the total number of feedback requests that the
  user responded to incorrectly
\end{itemize}

\subsubsection{User Errors per Class}\label{user-errors-per-class}

This table displays the details of user errors \textbf{per class} for:

\begin{itemize}
\tightlist
\item
  seed stage
\item
  review stage
\item
  in total
\end{itemize}

\subsection{Image Classification
Validation}\label{image-classification-validation}

The final section displays every image in your dataset with its Ground
Truth and Label Propagation classifications side by side. This detailed
view enables you to:

\begin{itemize}
\tightlist
\item
  Review individual classification decisions
\item
  Filter images by Ground Truth class, Label Propagation class, match
  status, or label source
\item
  Identify patterns in misclassifications
\item
  Audit specific defect types requiring attention
\end{itemize}

\begin{figure}
\centering
\pandocbounded{\includegraphics[keepaspectratio,alt={Image Classification Validation Example}]{images/vl-camtek-evaluation-imageclassification-example.png}}
\caption{Image Classification Validation Example}
\end{figure}

\processstart

To manipulate data in this table: \processend 

\begin{enumerate}
\def\labelenumi{\arabic{enumi}.}
\item
  Sort by any column. Click the title of the column to change sorting
  direction.
\item
  Filter with these options:

  \begin{itemize}
  \tightlist
  \item
    GT Classification - view images by a specific ground truth
    classification
  \item
    LP Classification - view images by a specific label propagation
    classification
  \item
    Result - filter by whether the label propagation was correct or
    incorrect
  \item
    Source - filter by the source of the label application - Visual
    Layer (the model), seed or the user (during review stages)
  \end{itemize}
\item
  View the continuation of the table from the page arrows at the top
  right of the table.
\item
  Click \textbf{Clear Filters} to view the entire table.
\end{enumerate}

\chapter{Export Results}\label{export-results}

Export detailed labeling results for quality analysis, auditing, or
integration with external tools. Export is available at any stage of
label propagation, even before reaching train-worthy status. Use this
for quality audits or to track annotation progress across iterations.

\notestart

This export is optional and not required for model training through the
integrated workflow. \noteend

\section{Export Options}\label{export-options}

\textbf{Visual Layer} generates a JSON file that includes detailed
attribution for every labeled image, enabling you to track model
accuracy and annotation quality. You can export this file as follows:

\begin{longtable}[]{@{}
  >{\raggedright\arraybackslash}p{(\linewidth - 2\tabcolsep) * \real{0.5000}}
  >{\raggedright\arraybackslash}p{(\linewidth - 2\tabcolsep) * \real{0.5000}}@{}}
\toprule\noalign{}
\begin{minipage}[b]{\linewidth}\raggedright
\textbf{Export Option}
\end{minipage} & \begin{minipage}[b]{\linewidth}\raggedright
\textbf{Description}
\end{minipage} \\
\midrule\noalign{}
\endhead
\bottomrule\noalign{}
\endlastfoot
\textbf{Export items that match the current filter selection} & Export
the JSON with only the metadata for a portion of the dataset \\
\textbf{Entire Dataset} & Export the JSON for the complete dataset \\
\textbf{Label Propagation Results} & Export the JSON and all label
results; this option only appears once you've completed label
propagation. \\
\textbf{Include image files} & Include up to 5,000 images or only export
the JSON \\
\end{longtable}

\processstart

\textbf{To export results} \processend

\begin{enumerate}
\def\labelenumi{\arabic{enumi}.}
\item
  Navigate to your dataset.
\item
  If relevant, filter the dataset to the results that you wish to
  export.
\item
  Click
  
\includegraphics[width=0.05\linewidth,height=\textheight,keepaspectratio,alt={Download icon}]{images/vl-export-download-icon.png}
  from the top right. The \textbf{Export data} options open.

  \begin{center}
  \setlength{\fboxrule}{5pt}
  \setlength{\fboxsep}{0pt}
  \fcolorbox{imageborder}{white}{\includegraphics[width=0.5\textwidth]{images/vl-export-options-dialog.png}}
  \end{center}
\item
  Configure the export with the relevant
  \hyperref[export-options]{options}.
\item
  Click \textbf{Export File}.

  The JSON file downloads automatically to your default downloads
  folder. If you selected to export the images as well, they are
  downloaded in a ZIP file simultaneously.
\end{enumerate}

\subsection{Example Structure for Dataset
Export}\label{example-structure-for-dataset-export}

Use this export to work with your dataset, or portion thereof, and all
of its metadata.

The structured output appears similar to the following example:

\begin{lstlisting}
{
  "dataset_id": "def456",
  "timestamp": "2025-10-25T10:00:00Z",
  "total_annotations": 950,
  "total_cycles": 3,
  "annotations": [
    {
      "image_path": "images/wafer1/die_101.jpg",
      "label_name": "Surface Defect - Major",
      "source": "VL",
      "cycle_number": 3
    },
    {
      "image_path": "images/wafer1/die_102.jpg",
      "label_name": "Good Quality",
      "source": "User",
      "cycle_number": 2
    },
    {
      "image_path": "images/wafer2/die_205.jpg",
      "label_name": "Color Inconsistency",
      "source": "Seed",
      "cycle_number": 1
    },
    {
      "image_path": "images/wafer3/die_315.jpg",
      "label_name": "Structural Damage",
      "source": "VL",
      "cycle_number": 3
    }
  ]
}
\end{lstlisting}

If you export the entire dataset or a filtered portion of the dataset,
the JSON file contains:

\begin{longtable}[]{@{}
  >{\raggedright\arraybackslash}p{(\linewidth - 4\tabcolsep) * \real{0.3333}}
  >{\raggedright\arraybackslash}p{(\linewidth - 4\tabcolsep) * \real{0.3333}}
  >{\raggedright\arraybackslash}p{(\linewidth - 4\tabcolsep) * \real{0.3333}}@{}}
\toprule\noalign{}
\begin{minipage}[b]{\linewidth}\raggedright
\textbf{Field Key}
\end{minipage} & \begin{minipage}[b]{\linewidth}\raggedright
\textbf{Data Type}
\end{minipage} & \begin{minipage}[b]{\linewidth}\raggedright
\textbf{Description}
\end{minipage} \\
\midrule\noalign{}
\endhead
\bottomrule\noalign{}
\endlastfoot
\textbf{info} & \texttt{Object} & Top-level object containing metadata
about the export and dataset. \\
\textbf{schema\textbackslash version} & \texttt{String} & The version of
the export schema. \\
\textbf{dataset} & \texttt{String} & The human-readable name of the
dataset. \\
\textbf{description} & \texttt{String} & A brief description of the
export operation. \\
\textbf{dataset\textbackslash url} & \texttt{String} & A URL pointing to
the dataset within the \textbf{Visual Layer} application. \\
\textbf{export\textbackslash time} & \texttt{String\ (ISO\ 8601)} &
Timestamp indicating when the data export was initiated. \\
\textbf{dataset\textbackslash creation\textbackslash time} &
\texttt{String\ (ISO\ 8601)} & Timestamp indicating when the dataset was
originally created. \\
\textbf{exported\textbackslash by} & \texttt{String} & The user who
performed the data export. \\
\textbf{total\textbackslash media\textbackslash items} &
\texttt{Integer} & The total number of images/media items in the
dataset. \\
\textbf{media\textbackslash items} & \texttt{Array\ of\ Objects} & Array
containing details for each image in the export (similar to the results
array in the standard example, but with richer media details). \\
\textbf{media\textbackslash id} & \texttt{String} & Unique identifier
for the media item (image). \\
\textbf{media\textbackslash type} & \texttt{String} & The type of media
(e.g., ``image''). \\
\textbf{file\textbackslash name} & \texttt{String} & The original file
name of the image. \\
\textbf{file\textbackslash path} & \texttt{String} & The path to the
file within the original AOI output structure. \\
\textbf{file\textbackslash size} & \texttt{String} & The size of the
image file. \\
\textbf{uniqueness\textbackslash score} & \texttt{Number\ (Float)} & A
score indicating how unique the image is compared to the rest of the
dataset. \\
\textbf{height} & \texttt{Integer} & The height of the image in
pixels. \\
\textbf{width} & \texttt{Integer} & The width of the image in pixels. \\
\textbf{url} & \texttt{String} & A URL to view the specific media item
in \textbf{Visual Layer}. \\
\textbf{cluster\textbackslash id} & \texttt{String} & The ID of the
visual cluster the image belongs to (from the Exploration features). \\
\textbf{metadata\textbackslash items} & \texttt{Array\ of\ Objects} &
Array containing application-specific or custom metadata related to the
image. \\
\textbf{type} & \texttt{String} & The type of metadata entry (e.g.,
``issue'', ``image\textbackslash label''). \\
\textbf{properties} & \texttt{Object} & Object containing specific
details for the metadata item type. \\
\textbf{issue\textbackslash type} & \texttt{String} & Specific
classification of the issue (e.g., ``mislabels''). \\
\textbf{issues\textbackslash description} & \texttt{String} & A
description or source of the issue. \\
\textbf{confidence} & \texttt{Number\ (Float)} & A confidence score
associated with the metadata/issue. \\
\end{longtable}

\subsection{Example Structure for Label Propagation
Export}\label{example-structure-for-label-propagation-export}

Use this export to track the source of each label to measure model
accuracy and identify areas for improvement:

\begin{longtable}[]{@{}
  >{\raggedright\arraybackslash}p{(\linewidth - 2\tabcolsep) * \real{0.5000}}
  >{\raggedright\arraybackslash}p{(\linewidth - 2\tabcolsep) * \real{0.5000}}@{}}
\toprule\noalign{}
\begin{minipage}[b]{\linewidth}\raggedright
\textbf{Use Case}
\end{minipage} & \begin{minipage}[b]{\linewidth}\raggedright
\textbf{Description}
\end{minipage} \\
\midrule\noalign{}
\endhead
\bottomrule\noalign{}
\endlastfoot
\textbf{Model acceptance rate} & How often VL automatic labels were
approved versus corrected \\
\textbf{Class-specific performance} & Which defect types the model
handles well versus those needing more training data \\
\textbf{Quality assurance} & Full audit trail for compliance and process
validation \\
\textbf{Training integration} & Traceable labeled data for your model
pipeline \\
\end{longtable}

The structured output appears similar to the following example:

\begin{lstlisting}
{
  "dataset_id": "abc123",
  "results": [
    {
      "image_id": "img_001",
      "label": "Surface Defect - Minor",
      "source": "VL",
      "confidence": 0.94
    },
    {
      "image_id": "img_002",
      "label": "Good Quality",
      "source": "User",
      "confidence": null
    },
    {
      "image_id": "img_003",
      "label": "Structural Damage",
      "source": "Seed",
      "confidence": null
    }
  ]
}
\end{lstlisting}

The JSON file contains the following information:

\begin{longtable}[]{@{}
  >{\raggedright\arraybackslash}p{(\linewidth - 4\tabcolsep) * \real{0.3333}}
  >{\raggedright\arraybackslash}p{(\linewidth - 4\tabcolsep) * \real{0.3333}}
  >{\raggedright\arraybackslash}p{(\linewidth - 4\tabcolsep) * \real{0.3333}}@{}}
\toprule\noalign{}
\begin{minipage}[b]{\linewidth}\raggedright
\textbf{Field}
\end{minipage} & \begin{minipage}[b]{\linewidth}\raggedright
\textbf{Data Type}
\end{minipage} & \begin{minipage}[b]{\linewidth}\raggedright
\textbf{Description}
\end{minipage} \\
\midrule\noalign{}
\endhead
\bottomrule\noalign{}
\endlastfoot
\textbf{dataset\textbackslash id} & \texttt{String} & A unique
identifier for the dataset. \\
\textbf{timestamp} & \texttt{String\ (ISO\ 8601)} & The date and time
when the label propagation results were exported. \\
\textbf{total\textbackslash annotations} & \texttt{Integer} & The total
number of images that received a label (either automatic or manual) in
this dataset. \\
\textbf{total\textbackslash cycles} & \texttt{Integer} & The total
number of label propagation iterations/cycles completed on this
dataset. \\
\textbf{annotations} & \texttt{Array\ of\ Objects} & A list containing
the labeling result details for each image. \\
\textbf{image\textbackslash path} & \texttt{String} & The file path of
the image within the original directory structure of the dataset. \\
\textbf{label\textbackslash name} & \texttt{String} & The final
class/label assigned to the image. \\
\textbf{source} & \texttt{String} & The origin of the assigned label:
user (manual assignment/review), seed (initial user-provided example),
or VL (\textbf{Visual Layer} automated assignment). \\
\textbf{cycle\textbackslash number} & \texttt{Integer} & The label
propagation iteration number when the label was assigned or last
reviewed. \\
\end{longtable}

\chapter{Task Manager}\label{task-manager}

The \textbf{Task Manager} provides a unified view of all dataset-related
operations across your \textbf{Visual Layer} workspace. It tracks
processes such as dataset creation, label propagation, and model
training, displaying their current status, progress, and any errors that
occurred.

Use the \textbf{Task Manager} to:

\begin{itemize}
\tightlist
\item
  Monitor active operations in real-time
\item
  Track operation history from the past 14 days
\item
  Identify and troubleshoot failed processes
\end{itemize}

\notestart

The \textbf{Task Manager} maintains a visible history of all operations
performed in the past 14 days. \noteend

The \textbf{Task Manager} displays near real-time updates on all
operations.

To access the \textbf{Task Manager}:

\begin{enumerate}
\def\labelenumi{\arabic{enumi}.}
\item
  Go to the \textbf{Navigation Panel} on the left side of the screen.
\item
  Click
  
\includegraphics[width=0.04\linewidth,height=\textheight,keepaspectratio,alt={Task Manager icon}]{images/vl-task-manager-icon.png}.

  The \textbf{Task Manager} opens, displaying all operations from the
  past 14 days, similar to the following:

  \pandocbounded{\includegraphics[keepaspectratio]{images/vl-task-manager.png}}
\end{enumerate}

\notestart

If no operations have been performed in the past 14 days, the table
appears empty with a notification indicating no recent activity.
\noteend

The \textbf{Task Manager} displays all operations in a comprehensive
table. Each row represents a single task and contains the following
information:

\rowcolors{1}{tableprimary}{tablesecondary}
\begin{tabularx}{\textwidth}{|p{0.25\textwidth}|X|}
\hline
\textbf{Column} & \textbf{Description} \\
\hline
\textbf{Dataset Name} & The name of the dataset where the operation is running or was performed. Click the dataset name to navigate directly to its Explore page. \\
\hline
\textbf{Dataset ID} & The unique identifier for the dataset. Use filters to search by specific dataset IDs. \\
\hline
\textbf{Job} & The job identifier from your manufacturing equipment metadata. Each dataset contains a single job value. Use the filter to view tasks by specific jobs. \\
\hline
\textbf{Setup} & The setup identifier from your manufacturing equipment metadata. Each dataset contains a single setup value. Use the filter to view tasks by specific setups. \\
\hline
\textbf{Recipe} & The recipe identifier(s) from your manufacturing equipment metadata. A dataset may contain multiple recipe values. If the recipe name exceeds the cell width, it appears truncated with "..." — hover over the cell to view all recipe values. Use the filter to view tasks by specific recipes. \\
\hline
\textbf{Task} & The type of operation being performed: \textbf{Dataset Creation}<!-- \textbf{Media Addition}, --> \textbf{Label Propagation}, \textbf{Training Model}, or \textbf{Re-indexing}. Use the filter to view specific task types. \\
\hline
\textbf{Start Time} & The timestamp when the operation was initiated. Use the filter to view tasks from the last 24 hours, last 7 days, or last 14 days. \\
\hline
\textbf{Duration} & The elapsed time since the operation started (for running tasks) or the total time the operation took (for completed tasks). Sort by this column to identify long-running operations. \\
\hline
\textbf{Status} & The current state of the operation. See the Status Types section below for detailed descriptions of each status. Use the filter to view tasks with specific statuses. \\
\hline
\textbf{Actions} & Retrieve Task ID.  \\
\hline
\end{tabularx}

\notestart

The \textbf{Task Manager} tracks operation-level status independently
from overall dataset status. A dataset can show \textbf{Ready} status in
the \textbf{Dataset Inventory} even when the \textbf{Task Manager}
displays completed tasks like \textbf{Training} since these operations
do not block dataset access. \noteend

\section{Status Types}\label{task-statuses}

All tasks display one of several statuses to indicate their current
state. Each status includes a percentage to show progress through the
operation.

\notestart

See \hyperref[task-manager-statuses]{Task Manager Status Types} in the
Appendices for detailed descriptions of all task statuses and their
meanings. \noteend

\section{Working with the Task Manager
Table}\label{working-with-the-task-manager-table}

The table displays tasks ordered by \textbf{Start Time} by default, with
the most recent operations appearing first. Each column in the
\textbf{Task Manager} supports filtering, grouping and sorting to help
you quickly locate specific operations:

\begin{itemize}
\tightlist
\item
  \textbf{Filtering}: Click the filter icon in any column header, select
  the desired values, and click \textbf{Apply} to update the table.
\item
  \textbf{Sorting}: Click the column header to sort tasks in ascending
  or descending order.
\end{itemize}

\appendixdivider{Appendices}

\chapter{User Interface Reference}\label{user-interface-reference}

This reference guide provides detailed descriptions of all
\textbf{Visual Layer} screens, components, and interface elements. Use
this section to understand what each field, button, and area of the
platform does.

\section{Home Page: Dataset Inventory}\label{dataset-inventory}

When you log in to \textbf{Visual Layer}, the home page displays the
\textbf{Dataset Inventory}:

\begin{figure}
\centering
\includegraphics[width=0.9\linewidth,height=\textheight,keepaspectratio,alt={Visual Layer Platform}]{images/vl-landing-overview.png}
\caption{\textbf{Visual Layer} Platform}
\end{figure}

The home page layout consists of:

\renewcommand{\arraystretch}{1.5}
\rowcolors{1}{tableprimary}{tablesecondary}

\begin{tabularx}{\textwidth}{|>{\centering\arraybackslash}m{4em}|p{0.35\textwidth}|X|}
\hline
\textbf{Area} & \textbf{Name} & \textbf{Description} \\
\hline

\raisebox{-7pt}{
\includegraphics[width=5.5mm]{images/1light.png}}
  & \textbf{Navigation Panel}
  & Navigate to: this area (the \textbf{Dataset Inventory}, which is your \textbf{Home} page), the \textbf{Task Manager}, the \textbf{Model Catalog}, and log in/log out \\
\hline

\raisebox{-7pt}{
\includegraphics[width=5.5mm]{images/2light.png}}
  & \textbf{Main Panel}
  & Search for datasets. \\
\hline

\raisebox{-7pt}{
\includegraphics[width=5.5mm]{images/3light.png}}
  & \textbf{Dataset Inventory}
  & This is your default Home page, as appears in the example above. The \textbf{Dataset Inventory} page displays a complete list of your datasets. Click any dataset to open it. \\
\hline

\raisebox{-6pt}{
\includegraphics[width=5.5mm]{images/4light.png}}
  & \multicolumn{1}{m{0.35\textwidth}|}{%
        
\includegraphics[width=4cm]{images/vl-create-dataset-create-button.png}%
    }
  & Create a new dataset. \\
\hline
\end{tabularx}

\subsection{Navigation Panel}\label{nav-panel}

The \textbf{Navigation Panel} provides access to all major areas of
\textbf{Visual Layer}:

\renewcommand{\arraystretch}{1.5}
\rowcolors{1}{tableprimary}{tablesecondary}

\begin{tabularx}{\textwidth}{|>{\centering\arraybackslash}m{4em}|p{0.35\textwidth}|X|}
\hline
\textbf{Icon} & \textbf{Name} & \textbf{Description} \\
\hline

\raisebox{-6pt}{\includegraphics[width=5.5mm]{images/vl-dataset-inventory-icon.png}}
  & \textbf{Dataset Inventory}
  & A complete list of your datasets. \\
\hline

\raisebox{-6pt}{
\includegraphics[width=5.5mm]{images/vl-task-manager-icon.png}}
  & \textbf{Task Manager}
  & Monitor and manage all active tasks, jobs, and processing workflows. \\
\hline

\raisebox{-6pt}{\includegraphics[width=5.5mm]{images/vl-modelcatalog-icon.png}}
  & \textbf{Model Catalog}
  & Browse, deploy, and manage trained models for inference and analysis. \\
\hline

\raisebox{-6pt}{
\includegraphics[width=5.5mm]{images/vl-avatar-icon.png}}
  & \textbf{Your Avatar}
  & Log out of your account. \\
\hline
\end{tabularx}

\subsection{Dataset Inventory View}\label{dataset-inventory-view}

The \textbf{Dataset Inventory} displays all datasets as cards:

\pandocbounded{\includegraphics[keepaspectratio]{images/vl-dataset-inventory-view.png}}

\pandocbounded{\includegraphics[keepaspectratio]{images/vl-dataset-inventory-cards.png}}

\subsubsection{Dataset Card}\label{dataset-card}

Each dataset card shows:

\begin{figure}[H]
\centering
\includegraphics[width=0.4\textwidth]{images/vl-dataset-inventory-card.png}
\caption{Dataset Card Example}
\end{figure}

\begin{longtable}[]{@{}
  >{\raggedright\arraybackslash}p{(\linewidth - 2\tabcolsep) * \real{0.5000}}
  >{\raggedright\arraybackslash}p{(\linewidth - 2\tabcolsep) * \real{0.5000}}@{}}
\toprule\noalign{}
\begin{minipage}[b]{\linewidth}\raggedright
\textbf{Field}
\end{minipage} & \begin{minipage}[b]{\linewidth}\raggedright
\textbf{Description}
\end{minipage} \\
\midrule\noalign{}
\endhead
\bottomrule\noalign{}
\endlastfoot
\textbf{Dataset Name} & The title of the dataset for easy identification
and reference. \\
\textbf{Creation Date} & The date when the dataset was created and added
to your inventory. \\
\textbf{Dataset Status} & Current state of the dataset: \textbf{Ready},
\textbf{Processing}, \textbf{Failed}, or \textbf{Archived}. \\
\textbf{Image Count} & Total number of images contained in the
dataset. \\
\textbf{Info Icon}

\includegraphics[width=0.03\linewidth,height=\textheight,keepaspectratio]{images/vl-info-i-icon.png}
& Hover to view associated Job Name, Setup Name, and Recipe Names. \\
\textbf{Three-Dot Menu} & Access dataset actions including
\textbf{Delete}. \\
\end{longtable}

\begin{figure}[H]
\centering
\includegraphics[width=0.7\textwidth]{images/vl-camtek-datasetinventory-metadata.png}
\caption{Dataset Metadata Popup Example}
\end{figure}

\subsection{Main Panel Controls}\label{main-panel-controls}

The \textbf{Main Panel} provides dataset management tools:

\includegraphics[width=1.05\linewidth,height=\textheight,keepaspectratio]{images/vl-main-nav-panel.png}

\begin{longtable}[]{@{}
  >{\raggedright\arraybackslash}p{(\linewidth - 2\tabcolsep) * \real{0.5000}}
  >{\raggedright\arraybackslash}p{(\linewidth - 2\tabcolsep) * \real{0.5000}}@{}}
\toprule\noalign{}
\begin{minipage}[b]{\linewidth}\raggedright
\textbf{Control}
\end{minipage} & \begin{minipage}[b]{\linewidth}\raggedright
\textbf{Description}
\end{minipage} \\
\midrule\noalign{}
\endhead
\bottomrule\noalign{}
\endlastfoot
\textbf{Search Bar} & Locate datasets by name. \\
\textbf{Filter by Job/Recipe/Setup} & Filter datasets by production
metadata (Job Name, Setup Name, Recipe Name). \\
\textbf{Sort By} & Arrange datasets based on status, creation date, or
other attributes. \\
\textbf{Clear Selection} & Reset applied filters. \\
\end{longtable}

\section{Dataset View}\label{dataset-view}

When you open a dataset, the main dataset exploration interface loads:

\begin{figure}
\centering
\includegraphics[width=0.9\linewidth,height=\textheight,keepaspectratio,alt={Dataset View with Details}]{images/vl-overview-cluster-view.png}
\caption{Dataset View with Details}
\end{figure}

The dataset view consists of:

\renewcommand{\arraystretch}{1.5}
\rowcolors{1}{tableprimary}{tablesecondary}

\begin{tabularx}{\textwidth}{|>{\centering\arraybackslash}m{4em}|p{0.35\textwidth}|X|}
\hline
\textbf{Area} & \textbf{Component} & \textbf{Description} \\
\hline

\raisebox{-7pt}{
\includegraphics[width=5.5mm]{images/1light.png}} & \textbf{Tabs} & Navigate between Explore, Data and Views. \\
\hline

\raisebox{-7pt}{
\includegraphics[width=5.5mm]{images/2light.png}} & \textbf{Filter Panel} & Apply and combine search criteria. \\
\hline

\raisebox{-7pt}{
\includegraphics[width=5.5mm]{images/3light.png}} & \textbf{Action Bar} & Access operations like exporting. \\
\hline

\raisebox{-7pt}{
\includegraphics[width=5.5mm]{images/4light.png}} & \textbf{Details Sidebar} & View metadata and insights, and access label propagation. \\
\hline

\raisebox{-7pt}{
\includegraphics[width=5.5mm]{images/5light.png}} & \textbf{Thumbnail Grid} & Visual representation of your data in clusters. \\
\hline
\end{tabularx}

\subsection{Dataset Metadata}\label{dataset-metadata}

Hover over the info icon

\includegraphics[width=0.03\linewidth,height=\textheight,keepaspectratio]{images/vl-info-i-icon.png}
next to the dataset name to view production metadata:

\begin{figure}[H]
\centering
\includegraphics[width=0.6\textwidth]{images/vl-camtek-dataset-metadata.png}
\caption{Dataset Metadata Popup}
\end{figure}

\begin{longtable}[]{@{}
  >{\raggedright\arraybackslash}p{(\linewidth - 2\tabcolsep) * \real{0.5000}}
  >{\raggedright\arraybackslash}p{(\linewidth - 2\tabcolsep) * \real{0.5000}}@{}}
\toprule\noalign{}
\begin{minipage}[b]{\linewidth}\raggedright
\textbf{Field}
\end{minipage} & \begin{minipage}[b]{\linewidth}\raggedright
\textbf{Description}
\end{minipage} \\
\midrule\noalign{}
\endhead
\bottomrule\noalign{}
\endlastfoot
\textbf{Job Name} & The production job identifier from the folder
metadata. \\
\textbf{Setup Name} & The setup configuration identifier from the folder
metadata. \\
\textbf{Recipe Names} & The recipe name(s) used during inspection. \\
\end{longtable}

\subsection{Cluster View}\label{clustering-granularity}

\textbf{Visual Layer} automatically groups similar images into clusters
during dataset creation. Each cluster is displayed as a preview card:

\begin{figure}
\centering
\includegraphics[width=0.93\linewidth,height=\textheight,keepaspectratio,alt={Cluster preview cards}]{images/vl-clusters-image-view.png}
\caption{Cluster preview cards}
\end{figure}

\begin{longtable}[]{@{}
  >{\raggedright\arraybackslash}p{(\linewidth - 2\tabcolsep) * \real{0.5000}}
  >{\raggedright\arraybackslash}p{(\linewidth - 2\tabcolsep) * \real{0.5000}}@{}}
\toprule\noalign{}
\begin{minipage}[b]{\linewidth}\raggedright
\textbf{Cluster Card Element}
\end{minipage} & \begin{minipage}[b]{\linewidth}\raggedright
\textbf{Description}
\end{minipage} \\
\midrule\noalign{}
\endhead
\bottomrule\noalign{}
\endlastfoot
\textbf{Number of items} & The total number of items that match your
current filters \\
\textbf{Prominent labels} & The most frequent labels within the
cluster \\
\textbf{Representing images} & Thumbnails of a few key example images
that visually summarize the cluster \\
\textbf{Visual Similarity Icon}

\includegraphics[width=0.07\linewidth,height=\textheight,keepaspectratio]{images/vl-visual-similarity-icon.png}
& Hover over a cluster to access visual similarity search \\
\end{longtable}

\subsection{Filter Menu}\label{filter-menu}

Access filters by clicking

\includegraphics[width=0.06\linewidth,height=\textheight,keepaspectratio]{images/vl-filterbutton.png}
in the dataset toolbar:

\begin{figure}
\centering
\includegraphics[width=0.5\linewidth,height=\textheight,keepaspectratio,alt={All Filters}]{images/vl-camtek-allfilters.png}
\caption{All Filters}
\end{figure}

\subsubsection{Available Filters}\label{available-filters}

\begin{longtable}[]{@{}
  >{\raggedright\arraybackslash}p{(\linewidth - 4\tabcolsep) * \real{0.3333}}
  >{\raggedright\arraybackslash}p{(\linewidth - 4\tabcolsep) * \real{0.3333}}
  >{\raggedright\arraybackslash}p{(\linewidth - 4\tabcolsep) * \real{0.3333}}@{}}
\toprule\noalign{}
\begin{minipage}[b]{\linewidth}\raggedright
\textbf{Filter Category}
\end{minipage} & \begin{minipage}[b]{\linewidth}\raggedright
\textbf{Description}
\end{minipage} & \begin{minipage}[b]{\linewidth}\raggedright
\textbf{Configuration}
\end{minipage} \\
\midrule\noalign{}
\endhead
\bottomrule\noalign{}
\endlastfoot
\textbf{Folders} & Filter by specific AOI output folders within your
dataset & Check the folders to include \\
\textbf{Files} & Filter by specific image files or file patterns & Enter
filename or pattern to match \\
\textbf{Labels} & Filter by assigned label classes & Select which
classes to include or exclude \\
\textbf{User Tags} & Filter by custom tags you've applied to images &
Select tags from your tag list \\
\textbf{Duplicates} & Show only duplicate or near-duplicate images &
Toggle on to show duplicates only \\
\textbf{Outliers} & Show images that differ significantly from typical
patterns & Toggle on to show outliers only \\
\textbf{Quality Issues} & Filter images flagged with quality problems &
Select issue types (blur, overexposure, underexposure, etc.) \\
\textbf{Select Uniques} & Show only unique images (exclude duplicates) &
Toggle on to show unique images only \\
\textbf{Insertion Time} & Filter by when images were added to the
dataset & Set date ranges for upload time \\
\textbf{Media Status} & Filter by processing or indexing status & Select
status types (pending, processing, ready, error) \\
\end{longtable}

\subsubsection{Applied Filters Display}\label{applied-filters-display}

Applied filters appear in the \textbf{Query Panel} above the image grid:

\begin{figure}
\centering
\includegraphics[width=0.8\linewidth,height=\textheight,keepaspectratio,alt={Applied Filters}]{images/vl-applied-filters.png}
\caption{Applied Filters}
\end{figure}

\begin{longtable}[]{@{}
  >{\raggedright\arraybackslash}p{(\linewidth - 2\tabcolsep) * \real{0.5000}}
  >{\raggedright\arraybackslash}p{(\linewidth - 2\tabcolsep) * \real{0.5000}}@{}}
\toprule\noalign{}
\begin{minipage}[b]{\linewidth}\raggedright
\textbf{Control}
\end{minipage} & \begin{minipage}[b]{\linewidth}\raggedright
\textbf{Description}
\end{minipage} \\
\midrule\noalign{}
\endhead
\bottomrule\noalign{}
\endlastfoot
\textbf{Filter Chip} & Each applied filter shows as a chip with its
criteria \\
\textbf{X Button} & Click to remove an individual filter \\
\textbf{Clear All} & Remove all applied filters at once \\
\end{longtable}

\notestart

Multiple filters combine using \textbf{AND} logic. \noteend

\subsection{Visual Similarity Search}\label{visual-search}

\textbf{Visual Layer} enables similarity search using deep visual
embeddings:

\begin{longtable}[]{@{}
  >{\raggedright\arraybackslash}p{(\linewidth - 2\tabcolsep) * \real{0.5000}}
  >{\raggedright\arraybackslash}p{(\linewidth - 2\tabcolsep) * \real{0.5000}}@{}}
\toprule\noalign{}
\begin{minipage}[b]{\linewidth}\raggedright
\textbf{Search Method}
\end{minipage} & \begin{minipage}[b]{\linewidth}\raggedright
\textbf{Action}
\end{minipage} \\
\midrule\noalign{}
\endhead
\bottomrule\noalign{}
\endlastfoot
\textbf{From cluster or image in cluster view} & Hover over a cluster or
the image and click

\includegraphics[width=0.07\linewidth,height=\textheight,keepaspectratio]{images/vl-visual-similarity-icon.png}
at the bottom of the card. \\
\textbf{From the Image Details Page} & Click a single image to open its
details page, then click

\includegraphics[width=0.3\linewidth,height=\textheight,keepaspectratio]{images/vl-visual-similarity-BUTTON.png}
from the \textbf{Action Bar}. \\
\textbf{Using an External Image} & Click

\includegraphics[width=0.08\linewidth,height=\textheight,keepaspectratio]{images/vl-visual-similarity-searchbaricon.png}
near the search bar inside your dataset, upload an image from your
device, optionally crop it, and start the search. \\
\end{longtable}

\subsection{Views Tab}\label{views-tab}

The \textbf{Views} tab displays saved filter combinations and search
queries:

\begin{longtable}[]{@{}
  >{\raggedright\arraybackslash}p{(\linewidth - 2\tabcolsep) * \real{0.5000}}
  >{\raggedright\arraybackslash}p{(\linewidth - 2\tabcolsep) * \real{0.5000}}@{}}
\toprule\noalign{}
\begin{minipage}[b]{\linewidth}\raggedright
\textbf{Control}
\end{minipage} & \begin{minipage}[b]{\linewidth}\raggedright
\textbf{Description}
\end{minipage} \\
\midrule\noalign{}
\endhead
\bottomrule\noalign{}
\endlastfoot
\textbf{Save as View} Button &

\includegraphics[width=0.6\linewidth,height=\textheight,keepaspectratio,alt={Save as View button}]{images/vl-save-view-button.png}
Saves current filters and search criteria as a reusable view \\
\textbf{View Name} Dialog &
\includegraphics[width=0.5\linewidth,height=\textheight,keepaspectratio,alt={Name the View dialog}]{images/vl-save-view-dialog.png}
Enter a descriptive, unique name for the view \\
\textbf{View Card} & Click any saved view to load its filter
combination \\
\textbf{Three-Dot Menu} & Access view actions including
\textbf{Delete} \\
\textbf{Delete Confirmation} &
\includegraphics[width=0.6\linewidth,height=\textheight,keepaspectratio,alt={Confirm Deletion}]{images/vl-dataset-views-delete-confirmation.png} \\
\textbf{Delete Notification} &

\includegraphics[width=0.65\linewidth,height=\textheight,keepaspectratio,alt={Delete Notification}]{images/vl-dataset-views-delete-notification.png} \\
\end{longtable}

\notestart

Views are immutable. To modify a view, delete it and create a new one
with updated criteria. \noteend

\chapter{Dataset and Media Validation
Requirements}\label{validation-requirements}

\label{validation-requirements}

When creating datasets or adding media to existing datasets,
\textbf{Visual Layer} validates that all data meets structural and
compatibility requirements. This ensures data integrity across your
datasets.

\section{Required File Structure}\label{required-file-structure}

All AOI output folders must contain:

\begin{longtable}[]{@{}
  >{\raggedright\arraybackslash}p{(\linewidth - 2\tabcolsep) * \real{0.5000}}
  >{\raggedright\arraybackslash}p{(\linewidth - 2\tabcolsep) * \real{0.5000}}@{}}
\toprule\noalign{}
\begin{minipage}[b]{\linewidth}\raggedright
\textbf{File}
\end{minipage} & \begin{minipage}[b]{\linewidth}\raggedright
\textbf{Requirement}
\end{minipage} \\
\midrule\noalign{}
\endhead
\bottomrule\noalign{}
\endlastfoot
\textbf{db.csv} & Must contain ImageName column with values matching
actual image filenames in the folder \\
\textbf{metadata.json} & Contains Job Name, Setup Name, Recipe Name, and
Tool List \\
\textbf{classes.json} & Class taxonomy definition (if present in
original dataset, must match exactly when adding media) \\
\textbf{db\_export\_*.json} & Must follow naming pattern
db\_export\_{[}id{]}.json \\
\end{longtable}

\section{Metadata Compatibility
Requirements}\label{metadata-compatibility-requirements}

For dataset creation, all selected folders must share identical
production metadata. When adding media to existing datasets, new folders
must exactly match the original dataset metadata:

\begin{longtable}[]{@{}
  >{\raggedright\arraybackslash}p{(\linewidth - 2\tabcolsep) * \real{0.5000}}
  >{\raggedright\arraybackslash}p{(\linewidth - 2\tabcolsep) * \real{0.5000}}@{}}
\toprule\noalign{}
\begin{minipage}[b]{\linewidth}\raggedright
\textbf{Parameter}
\end{minipage} & \begin{minipage}[b]{\linewidth}\raggedright
\textbf{Requirement}
\end{minipage} \\
\midrule\noalign{}
\endhead
\bottomrule\noalign{}
\endlastfoot
\textbf{Job Name} & Must match across all folders (or match original
dataset when adding media) \\
\textbf{Setup Name} & Must match across all folders (or match original
dataset when adding media) \\
\textbf{Recipe Name} & Must match across all folders (or match original
dataset when adding media) \\
\textbf{Tool List} & Must match across all folders (or match original
dataset when adding media) \\
\textbf{Classes Taxonomy} & Must match when adding media to existing
datasets \\
\end{longtable}

\section{Common Validation Errors}\label{common-validation-errors}

If validation fails, \textbf{Visual Layer} displays specific error
messages:

\begin{longtable}[]{@{}
  >{\raggedright\arraybackslash}p{(\linewidth - 4\tabcolsep) * \real{0.3333}}
  >{\raggedright\arraybackslash}p{(\linewidth - 4\tabcolsep) * \real{0.3333}}
  >{\raggedright\arraybackslash}p{(\linewidth - 4\tabcolsep) * \real{0.3333}}@{}}
\toprule\noalign{}
\begin{minipage}[b]{\linewidth}\raggedright
\textbf{Error}
\end{minipage} & \begin{minipage}[b]{\linewidth}\raggedright
\textbf{Cause}
\end{minipage} & \begin{minipage}[b]{\linewidth}\raggedright
\textbf{Solution}
\end{minipage} \\
\midrule\noalign{}
\endhead
\bottomrule\noalign{}
\endlastfoot
\textbf{Missing required files} & Incomplete folder structure or folder
lacks required files & Verify AOI machine completed the scan
successfully and all required files are present \\
\textbf{Invalid file naming} & db\_export file doesn't match pattern &
Check filename follows db\_export\_{[}id{]}.json format \\
\textbf{JSON parsing error} & Syntax errors in classes.json or
db\_export & Validate JSON formatting \\
\textbf{Cross-validation failure} & Job/Setup/Recipe mismatch between
files within a folder & Verify folder integrity; may need to regenerate
scan \\
\textbf{Job Recipe mismatch} & Job Name, Setup Name, Recipe Name, or
Tool List differs across selected folders or from original dataset &
Verify folders are from the same production run, or when adding media,
verify the metadata.json of the data folder matches the original dataset
exactly \\
\textbf{ImageName mismatch} & Filenames in db.csv don't match actual
image files & Check that the ImageName column in db.csv contains exact
filenames including extensions \\
\textbf{Class taxonomy mismatch} & classes.json defines different
classes than the original dataset & When adding media, use the same
classes.json file from your original dataset creation \\
\end{longtable}

\chapter{Dataset Creation Stages}\label{dataset-creation-stages}

\label{dataset-creation-stages}

During dataset creation, \textbf{Visual Layer} prepares the dataset for
exploration through the following stages:

\begin{longtable}[]{@{}
  >{\raggedright\arraybackslash}p{(\linewidth - 2\tabcolsep) * \real{0.5000}}
  >{\raggedright\arraybackslash}p{(\linewidth - 2\tabcolsep) * \real{0.5000}}@{}}
\toprule\noalign{}
\begin{minipage}[b]{\linewidth}\raggedright
\textbf{Status}
\end{minipage} & \begin{minipage}[b]{\linewidth}\raggedright
\textbf{Description}
\end{minipage} \\
\midrule\noalign{}
\endhead
\bottomrule\noalign{}
\endlastfoot
\textbf{NEW} & Dataset initialization started \\
\textbf{UPLOADING} & Upload request acknowledged, data transfer in
progress \\
\textbf{INITIALIZING} & Pre-processing setup in progress \\
\textbf{PRE\_PROCESSING} & Data prepared for indexing and analysis \\
\textbf{INDEXING} & Visual indexing and AI-based analysis in progress:
learning distribution patterns and creating enhanced embeddings tailored
to your data for label propagation accuracy \\
\textbf{SAVING} & Final data save to system in progress \\
\textbf{READY} & Dataset complete and available for exploration \\
\textbf{FATAL\_ERROR} & Irrecoverable failure during dataset creation;
check error details in notification or \textbf{Task Manager} \\
\end{longtable}

\chapter{Dataset Operation Statuses}\label{dataset-operation-statuses}

After dataset creation (or dataset update), dataset status appears on
its card in the \textbf{Dataset Inventory} and inside the dataset at the
top right. Statuses indicate the current state and availability of
dataset operations:

\begin{longtable}[]{@{}
  >{\raggedright\arraybackslash}p{(\linewidth - 2\tabcolsep) * \real{0.5000}}
  >{\raggedright\arraybackslash}p{(\linewidth - 2\tabcolsep) * \real{0.5000}}@{}}
\toprule\noalign{}
\begin{minipage}[b]{\linewidth}\raggedright
\textbf{Status}
\end{minipage} & \begin{minipage}[b]{\linewidth}\raggedright
\textbf{Description}
\end{minipage} \\
\midrule\noalign{}
\endhead
\bottomrule\noalign{}
\endlastfoot
\textbf{READY} & No blocking operations are running. All actions are
available. \\
\textbf{UPDATING} & A blocking operation is running (\textbf{Label
Propagation} or \textbf{Re-indexing}). Most dataset actions are
unavailable until the operation completes. \\
\textbf{CREATING} & The dataset is being created. No dataset actions are
available until creation completes. \\
\textbf{RESUME} & User action is required to complete dataset creation.
Data upload started but \textbf{Create Dataset} was not clicked. \\
\textbf{PENDING} & New media was added and requires re-indexing. \\
\end{longtable}

\notestart

A dataset can show \textbf{Ready} status in the \textbf{Dataset
Inventory} even when the \textbf{Task Manager} displays completed tasks
like \textbf{Training} since these operations do not block dataset
access. \noteend

\chapter{Task Manager Statuses}\label{task-manager-statuses}

\label{task-manager-statuses}

Tasks are operations you perform in the system, such as creating a
dataset, updating an existing dataset or running label propagation. All
tasks in the \textbf{Task Manager} display one of the following statuses
to indicate their current state. Each status includes a percentage to
show progress through the operation:

\begin{longtable}[]{@{}
  >{\raggedright\arraybackslash}p{(\linewidth - 2\tabcolsep) * \real{0.5000}}
  >{\raggedright\arraybackslash}p{(\linewidth - 2\tabcolsep) * \real{0.5000}}@{}}
\toprule\noalign{}
\begin{minipage}[b]{\linewidth}\raggedright
\textbf{Status}
\end{minipage} & \begin{minipage}[b]{\linewidth}\raggedright
\textbf{Description}
\end{minipage} \\
\midrule\noalign{}
\endhead
\bottomrule\noalign{}
\endlastfoot
\textbf{Initializing} & The operation is performing initial validations,
uploading data, or downloading resources. Progress percentage reflects
completion of initialization steps. \\
\textbf{Running} & The operation is actively processing. Progress
percentage reflects completion of the overall operation. \\
\textbf{Succeeded} & The operation completed successfully. Progress
shows 100\%. \\
\textbf{Failed} & The operation encountered an error and could not
complete. Hover over the status to view the error message explaining
what went wrong. \\
\textbf{Aborted} & The operation was manually stopped by a user before
completion. \\
\textbf{Waiting for Review} & The label propagation operation completed
and is ready for your review. Navigate to the dataset to begin reviewing
propagated labels. This status applies only to \textbf{Label
Propagation} tasks. \\
\end{longtable}

\chapter{Troubleshooting Common
Issues}\label{troubleshooting-common-issues}

\section{Dataset Creation}\label{dataset-creation}

\notestart

For detailed validation requirements and error messages, see
\hyperref[validation-requirements]{Dataset and Media Validation
Requirements}. \noteend

\begin{longtable}[]{@{}
  >{\raggedright\arraybackslash}p{(\linewidth - 4\tabcolsep) * \real{0.3333}}
  >{\raggedright\arraybackslash}p{(\linewidth - 4\tabcolsep) * \real{0.3333}}
  >{\raggedright\arraybackslash}p{(\linewidth - 4\tabcolsep) * \real{0.3333}}@{}}
\toprule\noalign{}
\begin{minipage}[b]{\linewidth}\raggedright
\textbf{Issue}
\end{minipage} & \begin{minipage}[b]{\linewidth}\raggedright
\textbf{What to Check}
\end{minipage} & \begin{minipage}[b]{\linewidth}\raggedright
\textbf{Action}
\end{minipage} \\
\midrule\noalign{}
\endhead
\bottomrule\noalign{}
\endlastfoot
Upload fails or folders won't combine & Folder structure and metadata
compatibility & Verify AOI scan completed successfully. Check that all
required files exist (db.csv, metadata.json, classes.json,
db\_export\_*.json). When combining folders or adding media, confirm
Job/Setup/Recipe/Tool List match exactly. \\
Validation errors about file naming or format & File structure issues &
Check db\_export files follow naming pattern. Validate JSON formatting
in classes.json. Verify ImageName column in db.csv matches actual
filenames. \\
Class taxonomy errors when adding media & Classes don't match original
dataset & Use the same classes.json from your original dataset creation.
Class definitions cannot change when adding media. If this error
persists, contact your system administrator. \\
\end{longtable}

\section{Label Propagation}\label{label-propagation}

\begin{longtable}[]{@{}
  >{\raggedright\arraybackslash}p{(\linewidth - 4\tabcolsep) * \real{0.3333}}
  >{\raggedright\arraybackslash}p{(\linewidth - 4\tabcolsep) * \real{0.3333}}
  >{\raggedright\arraybackslash}p{(\linewidth - 4\tabcolsep) * \real{0.3333}}@{}}
\toprule\noalign{}
\begin{minipage}[b]{\linewidth}\raggedright
\textbf{Problem}
\end{minipage} & \begin{minipage}[b]{\linewidth}\raggedright
\textbf{Description}
\end{minipage} & \begin{minipage}[b]{\linewidth}\raggedright
\textbf{Resolution}
\end{minipage} \\
\midrule\noalign{}
\endhead
\bottomrule\noalign{}
\endlastfoot
Inaccurate results for certain defect types & The model consistently
mislabels specific defect types. High-confidence automatic labels are
frequently incorrect during review. & Add more diverse seed examples
showing the full range of that defect type's visual variations. Avoid
using edge cases or ambiguous examples as initial seeds. Reset if needed
to start with improved seeds. \\
Cannot drag multiple selected items to assign labels & You're trying to
drag multiple selected images to a class label, but only single items
can be dragged. & Use the Select \& Add workflow instead: select
multiple images using checkboxes, click \textbf{Assign Label}, then
choose the target class from the dropdown. \\
\end{longtable}

\chapter{Frequently Asked Questions}\label{frequently-asked-questions}

\textbf{Q: Why do review batches keep appearing? How many times do I
need to review?}

A: Label propagation is an iterative process where each review round
improves the model's accuracy. Continue reviewing until Visual Layer
reaches high confidence across your dataset or you're satisfied with the
labeling progress. The number of iterations varies based on dataset
complexity and seed quality.

\textbf{Q: Why can't I edit my seed examples after starting label
propagation?}

A: Seeds become immutable once propagation begins to ensure consistency
throughout the iterative learning process. If you need to change seeds
or classes, use Reset to start over. Note that this permanently deletes
all label propagation progress but preserves your original dataset.

\textbf{Q: What do the source indicators (VL, User badges) mean?}

A: These badges show the origin of each label. \textbf{VL} means the
label was assigned by Visual Layer with high confidence. \textbf{User}
means you reviewed and confirmed the label in a previous iteration. Seed
examples appear without a badge. See
\hyperref[reviewing-and-providing-feedback]{Reviewing and Providing
Feedback} for complete definitions.

\textbf{Q: Can I run label propagation on the same dataset multiple
times?}

A: Yes. However, each new run will overwrite the previous results within
that dataset.

\textbf{Q: What happens to images I marked as ``Ignore''?}

A: They are excluded from label propagation, placed in an ``Unlabeled''
cluster, and don't count toward training readiness or review quotas.

\textbf{Q: Can I export data before reaching ``Train Worthy'' status?}

A: Yes. You can export label propagation results at any point using
Export, even if you haven't reached the 100 images per class threshold.

\textbf{Q: How do I know if my seed examples are good quality?}

A: Monitor first-iteration results. If many high-confidence automatic
labels are incorrect during your first review, your seed likely needs
improvement. Reset and select clearer, more representative examples.

\textbf{Q: Can I combine datasets from different Jobs/Setups/Recipes?}

A: No.~Each dataset must contain folders with identical Job, Setup, and
Recipe metadata. Create separate datasets for different production runs,
then filter and organize them using the metadata filters.

\textbf{Q: Why are some folders grayed out when I try to add media or
create a dataset?}

A: Visual Layer automatically validates metadata compatibility.
Grayed-out folders indicate their Job, Setup, Recipe, or Tool List
metadata doesn't match the other selected folders. Hover over the info
icon to view the specific mismatch. You can either select only folders
with matching metadata, or create separate datasets for incompatible
production runs.

\textbf{Q: How do I fix incorrect seed examples or classes after
starting label propagation?}

A: Use Reset to start over with new seeds or classes. This permanently
deletes all label propagation progress but preserves your original
dataset. Seeds become immutable once propagation begins to ensure
consistency throughout the iterative learning process, so resetting is
the only way to change them.

\textbf{Q: Why haven't I reached train-worthy status after multiple
review rounds?}

A: Visual Layer requires a minimum of 100 images per class for training
readiness. Continue review rounds until this threshold is met. If
progress is very slow, verify you have sufficient unlabeled data in your
dataset and consider adding more images. The number of iterations needed
varies based on dataset size, complexity, and seed quality.

\chapter{Terms and Definitions}\label{terms-and-definitions}

\begin{longtable}[]{@{}
  >{\raggedright\arraybackslash}p{(\linewidth - 4\tabcolsep) * \real{0.3333}}
  >{\raggedright\arraybackslash}p{(\linewidth - 4\tabcolsep) * \real{0.3333}}
  >{\raggedright\arraybackslash}p{(\linewidth - 4\tabcolsep) * \real{0.3333}}@{}}
\toprule\noalign{}
\begin{minipage}[b]{\linewidth}\raggedright
\textbf{Term}
\end{minipage} & \begin{minipage}[b]{\linewidth}\raggedright
\textbf{Definition}
\end{minipage} & \begin{minipage}[b]{\linewidth}\raggedright
\textbf{Example/Context}
\end{minipage} \\
\midrule\noalign{}
\endhead
\bottomrule\noalign{}
\endlastfoot
\textbf{Class} & Each distinct type of label in your classification
system. Classes are the buckets that organize your labels. You need at
least two classes for any classification task, and each class needs
enough examples to teach the system what belongs in that category. & If
classifying defects, your classes might be ``Scratch,''
``Contamination,'' ``Void,'' and ``Pass'' \\
\textbf{Class Taxonomy} & The complete hierarchical structure of all
possible classes in your classification system, defined in classes.json.
This cannot be changed after dataset creation. & Your defect taxonomy
might include top-level classes like ``Physical Defects'' and
``Contamination,'' each with specific sub-types \\
\textbf{Confidence Score} & A numerical measure (0-100\%) indicating how
certain the system is about an automatically assigned label. Higher
scores mean greater certainty. & An image labeled ``Scratch'' with 95\%
confidence versus one with 60\% confidence \\
\textbf{Dataset} & A collection of images and their associated metadata
organized for analysis and labeling. A dataset can contain thousands of
images from one or more AOI output directories. & All inspection images
from a production run of Product X across multiple scans \\
\textbf{Embedding} & The technology that enables Visual Layer to find
visually similar images and automatically label thousands from just a
few examples, dramatically reducing manual labeling work. & Visual Layer
analyzes your images to understand visual patterns, allowing it to
recognize that Image A looks more like your ``Scratch'' seeds than your
``Void'' seeds \\
\textbf{Ground Truth} & Verified, accurate labels that serve as the
authoritative reference for training and evaluating machine learning
models. Ground truth is established through expert review and
validation, providing the foundation for measuring model accuracy and
performance. In Visual Layer, your seed examples and reviewed labels
form the ground truth for label propagation. & A set of 100 manually
verified defect images where domain experts have confirmed each label,
used to measure how accurately your model classifies new images \\
\textbf{Iteration} & A complete cycle of seed selection, automatic
labeling, review, and refinement. Running multiple iterations improves
labeling accuracy over time. & First iteration: provide seeds and review
results. Second iteration: adjust seeds based on results and review
again \\
\textbf{Job/Production Metadata} & The production configuration
parameters that define how inspection was performed. All folders in a
dataset must share the same Job, Setup, and Recipe to ensure consistent
scanning conditions. & Job: ``Product-X-2025'', Setup:
``Standard-Inspection'', Recipe: ``High-Resolution-Scan'' \\
\textbf{Label} & The category or classification assigned to a single
image. Think of it as the answer to ``what type is this?'' for one
specific data point. A label is an instance of a class. & This
particular image is labeled ``Scratch'' (the label is the specific
assignment to this one image) \\
\textbf{Label Propagation} & A semi-supervised learning process that
combines automated pattern recognition with human oversight to label
large datasets efficiently. The system learns from your seed examples,
automatically assigns labels based on confidence scores, and generates
review batches of uncertain images for your verification. Each iteration
refines the model as you approve or correct labels, creating a feedback
loop that improves accuracy until the system can label independently.
Each dataset maintains its own isolated model to prevent
cross-contamination. & After providing 10 seed examples of ``Void''
defects, the system automatically labels 800 similar images with high
confidence, sends 100 uncertain images for your review, then uses your
feedback to relabel and improve accuracy in the next iteration \\
\textbf{Metadata} & Descriptive information about your images and
inspection process, stored in db.csv and db\_export\_*.json files.
Includes production parameters, timestamps, and file locations. &
Information like scan date, machine ID, product type, and image file
paths \\
\textbf{Seed} & The initial labeled examples you provide for each class.
These training examples show the system what each class looks like in
practice. Seeds are the foundation that makes automatic labeling
possible. Without good seed examples, the system won't know what
patterns to look for. & Selecting 5-10 clear images of scratches to
teach the system what ``Scratch'' defects look like \\
\end{longtable}

\end{document}
